\documentclass[12pt, a4paper]{article}

% ── Idioma y codificación ──────────────────────────────────────────────────
\usepackage[utf8]{inputenc}
\usepackage[T1]{fontenc}
\usepackage[spanish, es-nodecimaldot]{babel}

% ── Matemáticas ───────────────────────────────────────────────────────────
\usepackage{amsmath, amssymb, amsthm}
\usepackage{mathtools}

% ── Fuente ────────────────────────────────────────────────────────────────
\usepackage{lmodern}
\usepackage{microtype}

% ── Geometría de página ───────────────────────────────────────────────────
\usepackage[
  top=2.5cm, bottom=2.5cm,
  left=2.8cm, right=2.8cm
]{geometry}

% ── Encabezado y pie ──────────────────────────────────────────────────────
\usepackage{fancyhdr}
\setlength{\headheight}{15pt}
\pagestyle{fancy}
\fancyhf{}
\fancyhead[L]{\small\nouppercase{\leftmark}}
\fancyhead[R]{\small\thepage}
\renewcommand{\headrulewidth}{0.4pt}

\fancypagestyle{plain}{%
  \fancyhf{}%
  \renewcommand{\headrulewidth}{0pt}%
}

% ── Figuras y diagramas ───────────────────────────────────────────────────
\usepackage{tikz}
\usetikzlibrary{arrows.meta, calc, babel}
\usepackage{float}
\usepackage{caption}
\usepackage{array}
\usepackage{tabularx}
\usepackage{booktabs}
\usepackage{setspace}
\usepackage{enumitem}

% ── Colores ───────────────────────────────────────────────────────────────
\usepackage{xcolor}
\definecolor{brandblue}{HTML}{000000}
\definecolor{linegray}{HTML}{000000}

\captionsetup{
  font=small,
  labelfont={bf,color=brandblue},
  textfont=it,
  justification=centering
}

\newcolumntype{Y}{>{\raggedright\arraybackslash}X}
\newcolumntype{C}{>{\centering\arraybackslash}X}
\setlength{\tabcolsep}{7pt}
\renewcommand{\arraystretch}{1.18}

% ── Hipervínculos ─────────────────────────────────────────────────────────
\usepackage[hidelinks,hypertexnames=false]{hyperref}

% ── Cajas para resultados destacados ──────────────────────────────────────
\usepackage{mdframed}
\usepackage[skins,breakable]{tcolorbox}
\mdfdefinestyle{resultado}{%
  linecolor=black,
  backgroundcolor=white,
  linewidth=0.8pt,
  innertopmargin=6pt,
  innerbottommargin=6pt,
  innerleftmargin=8pt,
  innerrightmargin=8pt,
  skipabove=8pt,
  skipbelow=8pt,
}

\tcbset{
  demostracionbox/.style={
    enhanced,
    breakable,
    frame hidden,
    interior hidden,
    borderline west={0.7pt}{0pt}{black},
    left=8pt, right=6pt,
    top=4pt, bottom=4pt,
    before skip=8pt, after skip=8pt,
  }
}

% ── Entornos matemáticos ──────────────────────────────────────────────────
\theoremstyle{plain}
\newtheorem{teorema}{Teorema}[section]
\newtheorem{proposicion}[teorema]{Proposición}
\newtheorem{lema}[teorema]{Lema}
\newtheorem{corolario}[teorema]{Corolario}

\theoremstyle{definition}
\newtheorem{definicion}[teorema]{Definición}
\newtheorem{ejemplo}{Ejemplo}[section]

\theoremstyle{remark}
\newtheorem*{nota}{Nota}
\newtheorem*{observacion}{Observación}

% ── Ecuaciones numeradas para referencias ────────────────────────────────
% Convierte todas las expresiones en display \[...\] en ecuaciones numeradas.
\let\olddisplaymath\displaymath
\let\oldenddisplaymath\enddisplaymath
\renewenvironment{displaymath}{\begin{equation}}{\end{equation}}

% ── Formato de secciones ──────────────────────────────────────────────────
\usepackage{titlesec}
\usepackage{etoolbox}
\titleformat{\section}
  {\large\bfseries\color{brandblue}}
  {\thesection.}{0.6em}{}
  [\vspace{2pt}\color{linegray}\titlerule\vspace{4pt}]

\titleformat{\subsection}
  {\normalsize\bfseries\color{brandblue}}
  {\thesubsection.}{0.5em}{}

\titleformat{\subsubsection}
  {\normalsize\bfseries\color{brandblue}}
  {\thesubsubsection.}{0.5em}{}

\titleformat{\paragraph}
  {\normalsize\bfseries\color{brandblue}}
  {}{0em}{}

\titlespacing*{\section}{0pt}{0pt}{0.75\baselineskip}
\titlespacing*{\subsection}{0pt}{0.95\baselineskip}{0.45\baselineskip}
\titlespacing*{\subsubsection}{0pt}{0.75\baselineskip}{0.35\baselineskip}
\titlespacing*{\paragraph}{0pt}{0.65\baselineskip}{0.25\baselineskip}

% ── Macros de conveniencia ────────────────────────────────────────────────
\newcommand{\R}{\mathbb{R}}
\newcommand{\mat}[1]{\mathbf{#1}}
\newcommand{\tr}{\operatorname{tr}}
\newcommand{\rg}{\operatorname{r}}
\newcommand{\diag}{\operatorname{diag}}
\newcommand{\Var}{\operatorname{Var}}
\newcommand{\Cov}{\operatorname{Cov}}
\newcommand{\argmax}{\operatorname*{arg\,max}}
\newcommand{\argmin}{\operatorname*{arg\,min}}
\newcounter{varsprimcount}
\newenvironment{variablesprimera}{%
  \par\smallskip
  \noindent\textit{Definición de variables (primera aparición): }
  \setcounter{varsprimcount}{0}%
}{%
  .\par\smallskip
}
\newcommand{\varitem}[2]{%
  \ifnum\value{varsprimcount}>0\relax; \fi
  $#1$ es #2%
  \stepcounter{varsprimcount}%
}
\newcommand{\iniciotema}[3]{%
  \newpage
  \begin{center}
    {\Large\bfseries Tema #1}\\[0.3cm]
    {\large #2}\\[0.2cm]
    \rule{0.6\linewidth}{0.6pt}
  \end{center}
  \setcounter{section}{0}
  \addtocontents{toc}{\protect\bigskip\protect\noindent\textbf{Tema #1: #3}\protect\medskip}
}

% ── Separación vertical y ritmo de lectura ────────────────────────────────
\raggedbottom
\setstretch{1.08}
\setlength{\parskip}{4pt}
\setlength{\parindent}{0pt}
\setlength{\abovedisplayskip}{7pt}
\setlength{\belowdisplayskip}{7pt}
\setlength{\abovedisplayshortskip}{6pt}
\setlength{\belowdisplayshortskip}{6pt}

% ── Listas con márgenes homogéneos ────────────────────────────────────────
\setlist[itemize]{leftmargin=1.35em, itemsep=2pt, topsep=4pt, parsep=0pt, partopsep=0pt}
\setlist[enumerate]{leftmargin=1.55em, itemsep=2pt, topsep=4pt, parsep=0pt, partopsep=0pt}
\setlist[description]{style=nextline, leftmargin=0pt, labelwidth=\linewidth, labelindent=0pt,
  labelsep=0pt, itemindent=0pt, listparindent=0pt, itemsep=4pt, topsep=4pt, parsep=0pt, partopsep=0pt}

% ── Demostraciones ────────────────────────────────────────────────────────
\renewcommand{\proofname}{Demostración}
\renewcommand{\qedsymbol}{$\square$}
\BeforeBeginEnvironment{proof}{\begin{tcolorbox}[demostracionbox]}
\AfterEndEnvironment{proof}{\end{tcolorbox}}
\apptocmd{\proof}{\small\itshape}{}{}

% ─────────────────────────────────────────────────────────────────────────
\begin{document}

% ══════════════════════════════════════════════════════════════════════════
%  PORTADA
% ══════════════════════════════════════════════════════════════════════════
\hypersetup{pageanchor=false}
\begin{titlepage}
  \thispagestyle{empty}
  \centering
  \vspace*{2.4cm}

  {\Large\scshape Grado en Estadística y Ciencia de Datos\par}
  \vspace{0.25cm}
  {\normalsize Curso 2025--2026\par}

  \vspace{1.35cm}
  {\Huge\bfseries Análisis Multivariante\par}
  \vspace{0.25cm}
  {\Large\bfseries Apuntes Unificados\par}
  \vspace{0.55cm}
  {\normalsize Teoría, demostraciones y formulario\par}

  \vspace{1.0cm}
  \rule{0.54\linewidth}{0.6pt}
  \vspace{0.95cm}

  \begin{minipage}{0.78\linewidth}
    \centering
    \normalsize
    Tema 1: Definición, tipos y formalización de datos\\
    Tema 2: Distribuciones multivariantes y dependencia\\
    Tema 3: Análisis de componentes principales\\
    Tema 4: Análisis discriminante
  \end{minipage}

  \vfill
  {\normalsize Luis López\par}
  \vspace{0.15cm}
  {\small Abril 2026\par}
\end{titlepage}
\hypersetup{pageanchor=true}

% ══════════════════════════════════════════════════════════════════════════
%  ÍNDICE
% ══════════════════════════════════════════════════════════════════════════
\tableofcontents

\section*{Convenciones de notación}
\addcontentsline{toc}{section}{Convenciones de notación}

Para mantener una exposición homogénea en todo el documento, utilizaremos las
siguientes convenciones formales, que se respetarán en todos los temas:
\begin{itemize}
  \item Los \textbf{escalares} se denotan en minúscula ($x,\lambda,r$); los
        \textbf{vectores}, en negrita minúscula
        ($\mathbf{x},\boldsymbol{\mu}$); y las \textbf{matrices}, en negrita
        mayúscula ($\mat{X},\mat{S},\mat{R}$).
  \item $\mat{X}\in\R^{n\times p}$ representa la matriz de datos con
        $n$ observaciones (filas) y $p$ variables (columnas); esta convención
        se mantendrá en todos los temas.
  \item El superíndice prima ($'$) indica trasposición y
        $\mat{A}^{-1}$ indica inversa cuando existe.
  \item $\tr(\mat{A})$, $|\mat{A}|$ y $\rg(\mat{A})$ denotan, respectivamente,
        traza, determinante y rango de $\mat{A}$.
  \item $E(\mathbf{X})$, $\Var(\mathbf{X})$ y $\Cov(\mathbf{X},\mathbf{Y})$
        denotan esperanza, varianza y covarianza poblacionales.
  \item En contexto muestral se usa barra superior para medias
        ($\bar{\mathbf{x}}$) y sombrero para estimadores
        ($\hat{\boldsymbol{\Sigma}}$, $\hat{\pi}_g$).
\end{itemize}

% ══════════════════════════════════════════════════════════════════════════
\section{Definición y tipos}
% ══════════════════════════════════════════════════════════════════════════

\subsection{¿Qué es el Análisis Multivariante?}

\begin{definicion}[Análisis Multivariante]
El análisis multivariante es el conjunto de métodos estadísticos diseñados
para estudiar \textbf{simultáneamente} varias variables observadas sobre una
misma unidad experimental u observacional, preservando su estructura conjunta
de dependencia.
\end{definicion}

En gran parte de las técnicas multivariantes, el objeto algebraico central es
una combinación lineal de variables observadas, cuyas ponderaciones se ajustan
según el objetivo inferencial (síntesis, discriminación, predicción o
descripción estructural). En este documento la denominaremos
\textbf{valor teórico}.
Esta construcción aparece de forma explícita en ACP, análisis factorial y
discriminante, entre otros métodos; no agota toda la disciplina, pero sí
proporciona un hilo algebraico común para una parte sustancial del análisis
multivariante.

\begin{mdframed}[style=resultado]
\[
  \text{Valor teórico} = w_1 X_1 + w_2 X_2 + \cdots + w_p X_p
\]
\end{mdframed}
\begin{variablesprimera}
  \varitem{X_i}{valor observado de la variable $i$}
  \varitem{w_i}{peso asignado a la variable $i$ por la técnica multivariante}
  \varitem{p}{número de variables incluidas en la combinación}
\end{variablesprimera}

El resultado es una variable escalar sintética que resume información en una
dirección del espacio de variables definida por el vector de pesos
$\mathbf{w}=(w_1,\dots,w_p)'$. Geométricamente, se trata de proyectar cada
observación sobre una dirección y analizar cuánto contenido estadístico retiene
esa proyección.

\subsection{Objetivos principales}

Desde una perspectiva aplicada, el análisis multivariante persigue cuatro metas
complementarias:

\begin{enumerate}
  \item \textbf{Resumir} el conjunto original en pocas variables derivadas,
        minimizando pérdida de información relevante.
        (\textit{ej.: componentes principales, análisis factorial}).

  \item \textbf{Detectar grupos} naturales en los datos cuando existan.
        (\textit{ej.: tipología de procesos productivos, de clientes, de
        partículas}).

  \item \textbf{Clasificar} nuevas observaciones en grupos previamente
        definidos.
        (\textit{ej.: clientes de un banco como fiables o no, personas como
        enfermas o no, partículas subatómicas como fermión o bosón}).

  \item \textbf{Relacionar} bloques de variables para estudiar estructura
        común y dependencia cruzada.
        (\textit{ej.: un conjunto de variables de capacidad intelectual y otro de
        resultados profesionales}).
\end{enumerate}

\subsection{Clasificación de las técnicas}

Existen diversas taxonomías. Seguimos la propuesta de \textbf{Hair, Tathan y
Black (1998)}, que organiza las técnicas según el \emph{tipo de relación entre
variables}. La lógica es distinguir tres marcos: dependencia (hay respuesta y
predictores), interdependencia (se analiza estructura interna sin respuesta
prefijada) y modelos no lineales.

% ─── Árbol de clasificación simplificado ─────────────────────────────────
\begin{figure}[H]
  \centering
  \begin{tikzpicture}[>=Stealth,
      every node/.style={font=\small, inner sep=3pt},
      level 1/.style={sibling distance=58mm, level distance=14mm},
      level 2/.style={sibling distance=28mm, level distance=13mm},
      edge from parent/.style={draw,->,>=Stealth}]
    \node[draw, rounded corners] {Técnicas multivariantes}
      child { node[draw, rounded corners, align=center] {Relación de\\dependencia} }
      child { node[draw, rounded corners, align=center] {Interdependencia} }
      child { node[draw, rounded corners, align=center] {Modelos\\no lineales} };
  \end{tikzpicture}
  \caption{Esquema general de clasificación (Hair et al., 1998).}
\end{figure}

\subsubsection{Relación de dependencia entre variables}

En este bloque se distingue explícitamente una variable respuesta (o un bloque
de respuestas) y un conjunto de variables explicativas. El objetivo es modelar
la distribución o el valor esperado de la respuesta condicionado a los
predictores. Es, por tanto, el marco natural cuando el problema estadístico es
de predicción o clasificación supervisada.

\paragraph{Una variable dependiente.}

\begin{description}
  \item[\textbf{Variable dependiente cuantitativa.}]
    \textbf{Regresión múltiple:} variación de una variable dependiente
    cuantitativa a partir de la variación de dos o más variables independientes.

  \item[\textbf{Variable dependiente cualitativa.}]
    \begin{itemize}
      \item \textbf{Análisis discriminante:} variable dependiente dicotómica o
            politómica; divide la muestra total en grupos basándose en una
            variable dependiente con clases conocidas. Sus objetivos son:
            \begin{itemize}
              \item Entender las diferencias entre grupos.
              \item Predecir la verosimilitud de que una entidad pertenezca a
                    una clase o grupo particular, basándose en varias variables
                    cuantitativas independientes.
            \end{itemize}
      \item \textbf{Regresión logística:} variable dependiente dicotómica con
            variables independientes cuantitativas.
      \item \textbf{Modelos logit:} variable dependiente dicotómica con variables
            independientes cuantitativas \emph{y} cualitativas.
    \end{itemize}
\end{description}

\paragraph{Varias variables dependientes.}

\begin{description}
  \item[\textbf{Variables dependientes cuantitativas en una sola relación.}]
    \begin{itemize}
      \item Con variables independientes cuantitativas:\\
            \textbf{Análisis de correlación canónica:} extensión de la regresión
            múltiple al caso de múltiples variables dependientes. Obtiene un
            conjunto de ponderaciones para las variables dependientes e
            independientes que proporciona la correlación máxima entre el
            conjunto de variables dependientes y el de independientes.
      \item Con variables independientes cualitativas:
            \begin{itemize}
              \item \textbf{Análisis de varianza multivariante (MANOVA).}
              \item \textbf{Análisis de correspondencias.}
            \end{itemize}
    \end{itemize}
  \item[\textbf{Variables dependientes cuantitativas en múltiples relaciones.}]
        \textbf{Modelos de ecuaciones estructurales.}
\end{description}

En síntesis, los métodos de dependencia se organizan por el tipo de variable
respuesta y por la cantidad de respuestas modeladas simultáneamente.

\subsubsection{Interdependencia entre variables}

En interdependencia no se fija a priori la dicotomía respuesta--explicativa; el
interés se desplaza a la geometría y estructura interna del conjunto de
variables. En este marco la interpretación geométrica (proyecciones,
distancias, agrupamientos) adquiere un papel central.

\begin{description}
  \item[\textbf{Estudio de la relación entre variables.}]
    \textbf{Análisis de componentes principales (ACP):} analiza las
    interrelaciones entre un gran número de variables y las explica en términos
    de sus dimensiones subyacentes comunes (\emph{factores}). Permite condensar
    la información de las variables originales en un conjunto más pequeño de
    variables con una pérdida mínima de información.

  \item[\textbf{Estudio de las relaciones entre casos.}]
    \textbf{Análisis de conglomerados (cluster):} agrupa observaciones
    similares para descubrir la estructura de grupos en los datos.

  \item[\textbf{Estudio de la estructura de un conjunto de objetos.}]
    Con variables cuantitativas: \textbf{Escalado multidimensional}, que
    transforma similitudes o disimilitudes de un conjunto de objetos en
    distancias representadas en un espacio multidimensional.
\end{description}

\subsubsection{Técnicas basadas en modelos no lineales}

\begin{itemize}
  \item \textbf{Método de superficie de respuesta:} estudio de las relaciones
        de dependencia entre las variables.
  \item \textbf{Análisis factorial no lineal:} relaciones de interdependencia.
\end{itemize}

Esta tercera familia recoge situaciones donde la relación entre variables no se
describe adecuadamente mediante estructura lineal global.

% ══════════════════════════════════════════════════════════════════════════
\section{Pasos del análisis}
% ══════════════════════════════════════════════════════════════════════════

Un flujo de trabajo multivariante riguroso puede organizarse en las diez etapas
siguientes:

\begin{enumerate}
  \item \textbf{Definición del problema de investigación} y objetivos del
        análisis.
  \item \textbf{Definición de las variables}, su tipo (cuantitativa o
        cualitativa) y tipo de relación que se establece entre ellas.
  \item \textbf{Selección de la técnica multivariante} en función del objetivo
        y de las variables estudiadas.
  \item \textbf{Determinación del tamaño muestral} y elementos necesarios para
        la recogida de datos.
  \item \textbf{Recogida de datos.}
  \item \textbf{Análisis previo de los datos}: detección de valores atípicos,
        datos perdidos, normalidad univariante.
  \item \textbf{Evaluación de los supuestos} de la técnica multivariante
        (normalidad multivariante, homocedasticidad, linealidad, etc.).
  \item \textbf{Estimación de los parámetros del modelo.}
  \item \textbf{Evaluación del ajuste del modelo.}
  \item \textbf{Interpretación de los resultados.}
\end{enumerate}

Aunque se presentan en orden secuencial, estas etapas deben entenderse como una
guía iterativa: en práctica, el diagnóstico de supuestos y la interpretación
suelen obligar a volver a etapas previas.

\begin{observacion}
Las etapas 6 y 7 son críticas: los métodos multivariantes son muy sensibles a
la presencia de outliers y al incumplimiento de los supuestos distribucionales.
Un análisis exploratorio previo riguroso ahorra errores de interpretación.
\end{observacion}

% ══════════════════════════════════════════════════════════════════════════
\section{Formalización de los datos}
% ══════════════════════════════════════════════════════════════════════════

\subsection{Variable univariante como vector}

Si disponemos de $n$ observaciones numéricas de una variable escalar $X$, cada
dato puede identificarse con la componente $i$-ésima de un \textbf{vector
columna} en $\R^n$. En términos formales, una muestra univariante es un
elemento de $\R^n$:
\[
  \mathbf{x} = \begin{pmatrix} x_1 \\ x_2 \\ \vdots \\ x_n \end{pmatrix}
  \in \R^n.
\]

\begin{ejemplo}
Edades de 3 personas:
\[
  \mathbf{x} = \begin{pmatrix} 18 \\ 21 \\ 34 \end{pmatrix}.
\]
\end{ejemplo}

\subsection{Variable vectorial y matriz de datos}

Cuando se miden $p$ variables sobre cada una de las $n$ unidades
observacionales, los datos se organizan naturalmente en una tabla de $n$ filas
y $p$ columnas. Cada columna recoge una variable y cada fila el perfil
multivariante completo de una unidad. Esta representación matricial permite
escribir de manera compacta medias, centrado, covarianzas, proyecciones y
transformaciones lineales.

\begin{definicion}[Matriz de datos]
Los valores de las $p$ variables escalares en cada uno de los $n$ elementos se
representan en la \textbf{matriz de datos} $\mat{X}$ de dimensión $n\times p$:
\begin{mdframed}[style=resultado]
\[
  \mat{X} = \begin{pmatrix}
    x_{11} & x_{12} & \cdots & x_{1p} \\
    x_{21} & x_{22} & \cdots & x_{2p} \\
    \vdots & \vdots & \ddots & \vdots \\
    x_{n1} & x_{n2} & \cdots & x_{np}
  \end{pmatrix}, \qquad
  \begin{array}{l}
    i = 1,\dots,n \quad \text{(individuo)} \\
    j = 1,\dots,p \quad \text{(variable)}
  \end{array}
\]
\end{mdframed}
\begin{variablesprimera}
  \varitem{n}{número de individuos u observaciones}
  \varitem{x_{ij}}{valor de la variable $j$ en el individuo $i$}
  \varitem{\mat{X}_{(j)}}{columna $j$ de la matriz de datos}
\end{variablesprimera}
También se escribe $\mat{X} = [\mat{X}_{(1)},\dots,\mat{X}_{(p)}] =
[\mathbf{x}_{ij}]$, donde $\mat{X}_{(j)}$ es la columna asociada a la variable
$j$ y la fila $i$ contiene el perfil multivariante de la observación $i$.
Esta convención (filas = individuos, columnas = variables) se mantendrá en todo
el documento para conservar una lectura algebraica uniforme.
\end{definicion}

% ══════════════════════════════════════════════════════════════════════════
\section{Elementos de álgebra lineal}
% ══════════════════════════════════════════════════════════════════════════

\subsection{Vectores}

\begin{definicion}[Vector traspuesto]
El vector traspuesto de $\mathbf{x}$, denotado $\mathbf{x}'$, es el vector
$\mathbf{x}$ con los mismos componentes dispuestos en fila:
\[
  \mathbf{x}' = (x_1, x_2, \dots, x_n).
\]
\end{definicion}

\begin{definicion}[Producto escalar o interno]
El producto escalar de dos vectores $\mathbf{x}$ e $\mathbf{y}$ de igual
dimensión es el número:
\[
  \mathbf{x}'\mathbf{y} = \sum_{i=1}^n x_i y_i.
\]
En particular,
\[
  \mathbf{x}'\mathbf{x}=\sum_{i=1}^n x_i^2\ge 0,
\]
con igualdad si y solo si $\mathbf{x}=\mathbf{0}$.
Esta definición fija simultáneamente una estructura algebraica (bilineal) y una
estructura geométrica (ángulos, normas y proyecciones).
\end{definicion}

\begin{observacion}
En análisis multivariante, el producto escalar fija la noción geométrica de
longitud, ángulo y proyección, y por tanto está en la base de distancias,
ortogonalidad y cambios de base.
\end{observacion}

\begin{definicion}[Norma de un vector]
La \textbf{norma} o longitud de un vector $\mathbf{x}$ es la raíz cuadrada de
su producto escalar consigo mismo:
\[
  \|\mathbf{x}\| = \sqrt{\mathbf{x}'\mathbf{x}} = \sqrt{x_1^2 + x_2^2 + \cdots + x_n^2}.
\]
Geométricamente, es la longitud del segmento que une el origen con el punto
$\mathbf{x}$.
\end{definicion}

\subsubsection{Relación entre producto escalar, normas y ángulo}

El producto escalar de dos vectores puede calcularse como el producto de sus
normas por el coseno del ángulo que forman:

\begin{mdframed}[style=resultado]
\[
  \mathbf{x}'\mathbf{y} = \|\mathbf{x}\|\,\|\mathbf{y}\|\,\cos\theta.
\]
\end{mdframed}

\begin{proposicion}[Desigualdad de Cauchy--Schwarz]
Para cualesquiera $\mathbf{x},\mathbf{y}\in\R^n$ se cumple:
\[
  |\mathbf{x}'\mathbf{y}| \le \|\mathbf{x}\|\,\|\mathbf{y}\|.
\]
La igualdad se alcanza si y solo si $\mathbf{x}$ y $\mathbf{y}$ son linealmente
dependientes.
\end{proposicion}

\begin{proof}
Para $t\in\R$, la cantidad
$\|\mathbf{x}-t\mathbf{y}\|^2=(\mathbf{x}-t\mathbf{y})'(\mathbf{x}-t\mathbf{y})$
es no negativa. Al desarrollar:
\[
  \|\mathbf{x}\|^2-2t\,\mathbf{x}'\mathbf{y}+t^2\|\mathbf{y}\|^2\ge 0.
\]
Este polinomio cuadrático en $t$ debe tener discriminante no positivo, luego
\[
  4(\mathbf{x}'\mathbf{y})^2-4\|\mathbf{x}\|^2\|\mathbf{y}\|^2\le 0
  \quad\Longrightarrow\quad
  |\mathbf{x}'\mathbf{y}| \le \|\mathbf{x}\|\,\|\mathbf{y}\|.
\]
La igualdad equivale a discriminante nulo, que ocurre exactamente cuando existe
$t$ tal que $\mathbf{x}=t\mathbf{y}$. Es decir, la igualdad caracteriza la
colinealidad entre ambos vectores.
\end{proof}

\begin{ejemplo}
Para $\mathbf{x} = \begin{pmatrix}a\\0\end{pmatrix}$ e
$\mathbf{y} = \begin{pmatrix}a\\c\end{pmatrix}$:
\[
  \|\mathbf{x}\| = a,\quad
  \|\mathbf{y}\| = \sqrt{a^2+c^2},\quad
  \cos\theta = \frac{a}{\sqrt{a^2+c^2}},
\]
\[
  \mathbf{x}'\mathbf{y} = a\cdot\sqrt{a^2+c^2}\cdot\frac{a}{\sqrt{a^2+c^2}} = a^2.
\]
\end{ejemplo}

\begin{observacion}
El producto escalar puede interpretarse como el producto de la norma de un
vector por la \emph{proyección} del otro sobre él.
\end{observacion}

\subsubsection{Vector unitario y ortogonalidad}

\begin{definicion}[Vector unitario o normalizado]
Un vector es \textbf{unitario} si su norma es 1:
\[
  \mathbf{u} = \frac{\mathbf{x}}{\|\mathbf{x}\|}, \qquad \mathbf{u}'\mathbf{u} = 1.
\]
\end{definicion}

\begin{definicion}[Ortogonalidad]
Dos vectores son \textbf{ortogonales} si su producto escalar es nulo:
\[
  \mathbf{x}'\mathbf{y} = \|\mathbf{x}\|\,\|\mathbf{y}\|\,\cos\theta = 0
  \quad\Longleftrightarrow\quad \theta = 90^\circ.
\]
\end{definicion}

\begin{definicion}[Ortonormalidad]
Dos vectores son \textbf{ortonormales} si son ortogonales \emph{y} ambos
tienen norma 1 (normalizados).
\end{definicion}

% ─────────────────────────────────────────────────────────────────────────
\subsection{Matrices}

\begin{definicion}[Tipos básicos de matrices]
En esta materia aparecen de forma recurrente varias familias de matrices. No
son una lista aislada: reaparecen de forma sistemática en covarianzas,
correlaciones, proyecciones y cambios de base:
\begin{itemize}
  \item \textbf{Cuadrada de orden $p$}: dimensión $p\times p$ (si $n\neq p$,
        la matriz es rectangular).
  \item \textbf{Traspuesta}: para $\mat{A}\in\R^{n\times p}$,
        $\mat{A}'\in\R^{p\times n}$ se obtiene intercambiando filas y columnas.
  \item \textbf{Traza}: $\tr(\mat{A})=\sum_i a_{ii}$, con
        $\tr(\mat{A})=\tr(\mat{A}')$.
  \item \textbf{Simétrica}: $\mat{A}=\mat{A}'$.
  \item \textbf{Nula}: todos sus elementos son cero.
  \item \textbf{Diagonal}: elementos no diagonales nulos.
  \item \textbf{Escalar}: diagonal con todos los valores diagonales iguales.
  \item \textbf{Identidad}: diagonal con unos en la diagonal principal.
\end{itemize}
\end{definicion}

\subsubsection{Rango de una matriz}

\begin{definicion}[Rango]
El \textbf{rango} de una matriz $\mat{A}$, denotado $\rg(\mat{A})$, es el
máximo número de columnas (equivalentemente, filas) linealmente independientes
de $\mat{A}$.
\end{definicion}

Las propiedades de rango más usadas en estos apuntes son:

\begin{mdframed}[style=resultado]
\begin{align}
  &0 \le \rg(\mat{A}) \le \min(n,p). \\
  &\text{Si }\mat{A}\text{ es cuadrada: }\rg(\mat{A}) \le p. \\
  &\text{Si }\mat{A}\text{ es regular de orden }p:\rg(\mat{A}) = p. \\
  &\text{Si }\mat{A}\text{ es singular de orden }p:\rg(\mat{A}) < p. \\
  &\rg(\mat{A}) = \rg(\mat{A}').\\
  &\rg(\mat{A}'\mat{A}) = \rg(\mat{A}\mat{A}') = \rg(\mat{A}) = \rg(\mat{A}').
\end{align}
\end{mdframed}

Estas identidades se usarán de forma recurrente al estudiar colinealidad,
redundancia de variables y dimensión efectiva de representación.

\begin{observacion}
El rango cuantifica la dimensión efectiva del subespacio generado por las
variables. En términos estadísticos, detectar rango deficiente equivale a
detectar redundancia lineal.
\end{observacion}

\begin{proposicion}[Carácter de invertibilidad en matrices cuadradas]
Sea $\mat{A}\in\R^{p\times p}$. Son equivalentes:
\begin{enumerate}
  \item $\mat{A}$ es invertible;
  \item $|\mat{A}|\neq 0$;
  \item $\rg(\mat{A})=p$;
  \item el sistema homogéneo $\mat{A}\mathbf{x}=\mathbf{0}$ solo tiene la
        solución trivial.
\end{enumerate}
\end{proposicion}

\begin{proof}
La equivalencia $(1)\Leftrightarrow(2)$ es el criterio clásico del
determinante. La equivalencia $(2)\Leftrightarrow(3)$ se sigue de que, para
matrices cuadradas, tener determinante no nulo equivale a tener rango máximo.
Finalmente, $(3)\Leftrightarrow(4)$ expresa que la nulidad es cero
(\emph{teorema rango--nulidad}), por lo que el núcleo contiene solo
$\mathbf{0}$. En consecuencia, invertibilidad, rango máximo y unicidad de
solución en el sistema homogéneo son tres formulaciones de un mismo hecho
estructural.
\end{proof}

\begin{observacion}
El rango indica la \emph{dimensión real} necesaria para representar el
conjunto de datos (número real de variables aleatorias de que disponemos).
Analizar el rango de una matriz de datos es clave para reducir el número de
variables sin pérdida de información.
\end{observacion}

% ─────────────────────────────────────────────────────────────────────────
\subsection{Operaciones con matrices}

Las operaciones matriciales no son meramente formales: constituyen el lenguaje
operativo de todo el análisis multivariante (centrado, covarianzas,
estandarización, proyecciones y reglas discriminantes).

\subsubsection{Suma de matrices}

Para sumar dos matrices es necesario que tengan el \textbf{mismo orden}:
\[
  \mat{A} + \mat{B} = [a_{ij} + b_{ij}], \quad i=1,\dots,n,\; j=1,\dots,p.
\]

Propiedades adicionales a la conmutativa y asociativa:
\[
  (\mat{A}+\mat{B})' = \mat{A}' + \mat{B}', \qquad
  \tr(\mat{A}+\mat{B}) = \tr(\mat{A}) + \tr(\mat{B}).
\]

\subsubsection{Multiplicación por un escalar}

\[
  k\mat{A} = \mat{A}k = [k\,a_{ij}].
\]

\subsubsection{Multiplicación de matrices}

Para multiplicar $\mat{A}$ y $\mat{B}$, el número de columnas de $\mat{A}$
debe coincidir con el número de filas de $\mat{B}$:

\begin{mdframed}[style=resultado]
\[
  \mat{A}_{n\times p} \times \mat{B}_{p\times m} = \mat{C}_{n\times m},
  \quad c_{ij} = \sum_{k=1}^p a_{ik}\,b_{kj}.
\]
\end{mdframed}

\begin{ejemplo}
\[
  \begin{pmatrix}4&2&3\\5&1&2\end{pmatrix}
  \times
  \begin{pmatrix}2&3\\1&5\\4&2\end{pmatrix}
  =
  \begin{pmatrix}22&28\\19&24\end{pmatrix}.
\]
\end{ejemplo}

\paragraph{Propiedades de la multiplicación de matrices.}

\begin{itemize}
  \item $\mat{A}(\mat{B}\mat{C}) = (\mat{A}\mat{B})\mat{C}$ \quad —asociativa—
  \item $\mat{A}\mat{B} \neq \mat{B}\mat{A}$ \quad —no conmutativa en general—
  \item $\mat{A}\mat{I} = \mat{I}\mat{A} = \mat{A}$
  \item $\mat{A}^2 = \mat{A}\mat{A}$ \quad (si $\mat{A}$ es cuadrada)
  \item $\tr(\mat{A}\mat{B}) = \tr(\mat{B}\mat{A})$ \quad (si $\mat{A}$ es cuadrada)
  \item $|\mat{A}\mat{B}| = |\mat{A}|\,|\mat{B}|$ \quad (si $\mat{A}$ y $\mat{B}$ son cuadradas del mismo orden)
  \item $(\mat{A}\mat{B}\mat{C})' = \mat{C}'\mat{B}'\mat{A}'$
\end{itemize}

\begin{observacion}
La no conmutatividad de la multiplicación es clave: el orden de las operaciones
codifica qué espacio se transforma (variables, observaciones o subespacios
proyectados).
\end{observacion}

% ─────────────────────────────────────────────────────────────────────────
\subsection{Determinantes}

\begin{definicion}[Determinante]
El determinante $|\mat{A}|$ es un \textbf{polinomio} de los elementos de la
matriz $\mat{A}$, es decir, un valor escalar calculado a partir de una matriz
cuadrada.
\end{definicion}

\begin{observacion}
En multivariante, el determinante de una matriz de covarianzas mide volumen
generalizado y aparece de forma natural en densidades gaussianas y en criterios
de dependencia global. En particular, $|\mat{S}|^{1/2}$ se interpreta como
escala volumétrica de la nube de datos tras incorporar la geometría de
covarianza.
\end{observacion}

\begin{ejemplo}
\[
  \mat{A} = \begin{pmatrix}3&2\\1&4\end{pmatrix} \Rightarrow
  |\mat{A}| = 3\cdot4 - 2\cdot1 = 10.
\]
\[
  \mat{B} = \begin{pmatrix}4&1&2\\2&5&1\\3&6&2\end{pmatrix} \Rightarrow
  |\mat{B}| = 9 \quad \text{(desarrollo por cofactores)}.
\]
\end{ejemplo}

\begin{mdframed}[style=resultado]
\begin{itemize}
  \item Si $|\mat{A}| = 0$: la matriz es \textbf{singular} (no invertible,
        dependencia lineal entre filas/columnas).
  \item Si $|\mat{A}| \neq 0$: la matriz es \textbf{regular} (invertible,
        filas y columnas linealmente independientes).
\end{itemize}
\end{mdframed}

% ─────────────────────────────────────────────────────────────────────────
\subsection{Inversa de una matriz}

\begin{definicion}[Matriz inversa]
La inversa de una matriz $\mat{A}$ es la matriz $\mat{A}^{-1}$ tal que
\[
  \mat{A}\mat{A}^{-1} = \mat{A}^{-1}\mat{A} = \mat{I}.
\]
Solo existe para matrices \textbf{cuadradas y regulares}.
\end{definicion}

Se calcula mediante la \textbf{matriz adjunta}:

\begin{mdframed}[style=resultado]
\[
  \mat{A}^{-1} = \frac{1}{|\mat{A}|}\,\mat{A}^\Lambda,
\]
donde $\mat{A}^\Lambda$ es la adjunta de $\mat{A}'$: matriz de adjuntos
$A_{ij} = (-1)^{i+j}|\mat{A}_{ij}|$, siendo $\mat{A}_{ij}$ la submatriz de
orden $n-1$ obtenida eliminando la fila $i$ y la columna $j$ de $\mat{A}'$.
\end{mdframed}

\begin{observacion}
La inversa de la covarianza (matriz de precisión) codifica dependencias
condicionales lineales y tiene papel central en inferencia gaussiana y
discriminación lineal. Por ello, su existencia y estabilidad numérica no son un
detalle técnico, sino una condición metodológica relevante.
\end{observacion}

\begin{ejemplo}
Para $\mat{A} = \begin{pmatrix}3&2\\1&4\end{pmatrix}$, $|\mat{A}|=10$:
\[
  \mat{A}^{-1} = \frac{1}{10}\begin{pmatrix}4&-2\\-1&3\end{pmatrix}
  = \begin{pmatrix}0{,}4&-0{,}2\\-0{,}1&0{,}3\end{pmatrix}.
\]
\end{ejemplo}

\begin{ejemplo}
Para $\mat{B} = \begin{pmatrix}4&1&2\\2&5&1\\3&6&2\end{pmatrix}$,
$|\mat{B}|=9$:
\[
  \mat{B}^{-1} = \frac{1}{9}\begin{pmatrix}4&10&-9\\-1&2&0\\-3&-21&18\end{pmatrix}
  \approx \begin{pmatrix}0{,}44&1{,}11&-1\\-0{,}11&0{,}22&0\\-0{,}33&-2{,}33&2\end{pmatrix}.
\]
\end{ejemplo}

% ─────────────────────────────────────────────────────────────────────────
\subsection{Autovalores y autovectores}

La lógica interna de este bloque es la siguiente: primero se define el problema
espectral ($|\mat{A}-\lambda\mat{I}|=0$), después se identifican direcciones
invariantes (autovectores), y finalmente se interpreta el resultado como cambio
de base adaptado a la transformación lineal.

\subsubsection{Autovalores (valores propios)}

\begin{definicion}[Autovalor]
Sea $\mat{A}$ una matriz cuadrada de orden $p$. Un escalar $\lambda$ es un
\textbf{autovalor} (valor propio o raíz latente) de $\mat{A}$ si:
\[
  |\mat{A} - \lambda\mat{I}| = 0.
\]
Esta es la \textbf{ecuación característica}: un polinomio de grado $p$ cuyas
raíces (con multiplicidad algebraica) son los autovalores de $\mat{A}$.
\end{definicion}

\begin{observacion}
La lectura conceptual es: la matriz define una transformación lineal y los
autovalores cuantifican los factores de escala en direcciones invariantes
autovectoriales.
\end{observacion}

\begin{ejemplo}
Calcular los autovalores de $\mat{A} = \begin{pmatrix}1&4\\9&1\end{pmatrix}$:
\begin{align}
  |\mat{A}-\lambda\mat{I}| = 0 \quad&\Longrightarrow\quad
  \begin{vmatrix}1-\lambda & 4\\9 & 1-\lambda\end{vmatrix} = 0 \\
  (1-\lambda)^2 - 36 = 0 \quad&\Longrightarrow\quad \lambda^2 - 2\lambda - 35 = 0 \\
  \lambda &= \frac{2\pm\sqrt{4+140}}{2} \quad\Longrightarrow\quad
  \lambda_1 = 7,\quad \lambda_2 = -5.
\end{align}
\end{ejemplo}

\paragraph{Propiedades de los autovalores.}

\begin{mdframed}[style=resultado]
\begin{itemize}
  \item $\displaystyle\sum_i \lambda_i = \tr(\mat{A})$.
  \item $\displaystyle\prod_i \lambda_i = |\mat{A}|$.
  \item Si $|\mat{A}| = 0$, al menos un $\lambda_i = 0$.
  \item $\rg(\mat{A})$ es el número de $\lambda_i$ distintos de 0.
  \item En una matriz diagonal, los autovalores son los elementos de la diagonal principal.
\end{itemize}
\end{mdframed}

Estas identidades conectan directamente álgebra lineal y estadística
multivariante: por ejemplo, en ACP los autovalores de $\mat{S}$ miden varianza
explicada.

\subsubsection{Autovectores (vectores propios)}

\begin{definicion}[Autovector]
Sea $\mat{A}$ cuadrada, $\lambda$ un autovalor y $\mathbf{x}\neq\mathbf{0}$
tal que:
\[
  \mat{A}\mathbf{x} = \lambda\mathbf{x}.
\]
$\mathbf{x}$ es un \textbf{autovector} (vector propio, característico o latente)
de $\mat{A}$ asociado a $\lambda$.
\end{definicion}

La ecuación $\mat{A}\mathbf{x} = \lambda\mathbf{x}$ equivale al sistema
homogéneo $(\mat{A}-\lambda\mat{I})\mathbf{x}=\mathbf{0}$. Por tanto, para un
$\lambda$ dado existe autovector no nulo si y solo si
$|\mat{A}-\lambda\mat{I}|=0$.

En general, la multiplicidad geométrica puede ser menor que la algebraica. En
el caso simétrico real, el teorema espectral garantiza una base ortonormal
completa de autovectores, hecho crucial para ACP y para la diagonalización de
formas cuadráticas.

\begin{ejemplo}
Para $\mat{A} = \begin{pmatrix}1&4\\9&1\end{pmatrix}$ con
$\lambda_1=7$, $\lambda_2=-5$:

\textit{Autovector asociado a $\lambda_1=7$:}
\[
  (\mat{A}-7\mat{I})\mathbf{x}_1 = \mathbf{0}
  \quad\Longrightarrow\quad
  \begin{pmatrix}-6&4\\9&-6\end{pmatrix}\begin{pmatrix}x_{11}\\x_{21}\end{pmatrix}=\mathbf{0}
  \quad\Longrightarrow\quad
  \mathbf{x}_1 = \begin{pmatrix}2\\3\end{pmatrix}.
\]

\textit{Autovector asociado a $\lambda_2=-5$:}
\[
  (\mat{A}+5\mat{I})\mathbf{x}_2 = \mathbf{0}
  \quad\Longrightarrow\quad
  \begin{pmatrix}6&4\\9&6\end{pmatrix}\begin{pmatrix}x_{12}\\x_{22}\end{pmatrix}=\mathbf{0}
  \quad\Longrightarrow\quad
  \mathbf{x}_2 = \begin{pmatrix}2\\-3\end{pmatrix}.
\]
\end{ejemplo}

\subsubsection{Autovectores normalizados}

\begin{definicion}[Autovector normalizado]
\[
  \mathbf{u}_i = \frac{\mathbf{x}_i}{\|\mathbf{x}_i\|}.
\]
\end{definicion}

\begin{ejemplo}
Tomando los autovectores del ejemplo anterior
$\mathbf{x}_1=\begin{pmatrix}2\\3\end{pmatrix}$ y
$\mathbf{x}_2=\begin{pmatrix}2\\-3\end{pmatrix}$, se tiene:
\[
  \|\mathbf{x}_1\|=\sqrt{2^2+3^2}=\sqrt{13},\qquad
  \|\mathbf{x}_2\|=\sqrt{2^2+(-3)^2}=\sqrt{13}.
\]
\begin{align}
  \mathbf{u}_1
  &= \frac{\mathbf{x}_1}{\|\mathbf{x}_1\|}
   = \begin{pmatrix}2/\sqrt{13}\\3/\sqrt{13}\end{pmatrix}
   \approx \begin{pmatrix}0{,}55\\0{,}83\end{pmatrix}, \\
  \mathbf{u}_2
  &= \frac{\mathbf{x}_2}{\|\mathbf{x}_2\|}
   = \begin{pmatrix}2/\sqrt{13}\\-3/\sqrt{13}\end{pmatrix}
   \approx \begin{pmatrix}0{,}55\\-0{,}83\end{pmatrix}.
\end{align}
Por tanto, la \textbf{matriz de autovectores normalizados} queda:
\[
  \mat{U}=\begin{pmatrix}0{,}55&0{,}55\\0{,}83&-0{,}83\end{pmatrix}.
\]
\end{ejemplo}

\paragraph{Propiedades de los autovectores.}

\begin{mdframed}[style=resultado]
\begin{itemize}
  \item Si $\mat{A}$ es simétrica, sus autovectores son \textbf{ortogonales}.
  \item Dada $\mat{A}$ simétrica, siendo $\mat{\Lambda}$ la matriz diagonal que
        contiene los autovalores y $\mat{U}$ la matriz de autovectores
        normalizados:
        \[
          \mat{A} = \mat{U}\mat{\Lambda}\mat{U}' \quad\Longrightarrow\quad
          \mat{A}^{-1} = \mat{U}\mat{\Lambda}^{-1}\mat{U}'.
        \]
        \[
          \mat{\Lambda} = \mat{U}'\mat{A}\mat{U}.
        \]
  \item Al aplicar una matriz cuadrada de orden $n$ a un vector de dimensión
        $n$, en general se transforma en dirección y magnitud.
  \item Los autovectores son aquellos vectores que, al ser transformados por la
        matriz, solo modifican su longitud (norma), no su posición en el espacio.
\end{itemize}
\end{mdframed}

% ─────────────────────────────────────────────────────────────────────────
\subsection{Formas cuadráticas}

\begin{definicion}[Forma cuadrática]
Las formas cuadráticas aparecen de forma sistemática en multivariante:
distancia de Mahalanobis, densidades gaussianas, criterios de optimización de
ACP y análisis discriminante.

Si transformamos un vector $\mathbf{x}$ mediante una transformación lineal
$\mathbf{y} = \mat{B}\mathbf{x}$, la norma al cuadrado del nuevo vector es:
\[
  \mathbf{y}'\mathbf{y} = \mathbf{x}'\mat{B}'\mat{B}\mathbf{x} = \mathbf{x}'\mat{A}\mathbf{x},
\]
donde $\mat{A} = \mat{B}'\mat{B}$ es cuadrada y simétrica.

La expresión $Q = \mathbf{x}'\mat{A}\mathbf{x}$ se denomina \textbf{forma
cuadrática} asociada a $\mat{A}$. Si $\mat{A}$ es simétrica, su signo queda
completamente determinado por el signo de sus autovalores.
\end{definicion}

\begin{observacion}
Desde el punto de vista geométrico, una forma cuadrática define niveles
elipsoidales (o degenerados) según la estructura espectral de $\mat{A}$; esta
lectura reaparece continuamente en covarianzas y en distancias tipo
Mahalanobis.
\end{observacion}

\subsubsection{Desarrollo explícito}

Para $\mat{A}_{2\times2}$:
\[
  Q = \mathbf{x}'\mat{A}\mathbf{x}
  = \begin{pmatrix}x_1&x_2\end{pmatrix}
    \begin{pmatrix}a_{11}&a_{12}\\a_{21}&a_{22}\end{pmatrix}
    \begin{pmatrix}x_1\\x_2\end{pmatrix}
  = a_{11}x_1^2 + (a_{12}+a_{21})x_1x_2 + a_{22}x_2^2.
\]

En general, para $\mat{A}_{p\times p}$:
\begin{mdframed}[style=resultado]
\[
  Q = \mathbf{x}'\mat{A}\mathbf{x} = \sum_i a_{ii}x_i^2 + \sum_{i<j}(a_{ij}+a_{ji})x_ix_j.
\]
\end{mdframed}

\subsubsection{Clasificación de formas cuadráticas}

\begin{mdframed}[style=resultado]
\begin{itemize}
  \item Para $\mathbf{x}=\mathbf{0}$: $Q=0$ siempre.
  \item \textbf{Definida positiva:} $Q>0$ para todo $\mathbf{x}\neq\mathbf{0}$.
        Equivale a: $\mat{A}$ regular, $\rg(\mat{A})=p$ y todos los autovalores
        son positivos.
  \item \textbf{Semidefinida positiva:} $Q\ge0$ para todo $\mathbf{x}\neq\mathbf{0}$.
        Equivale a: $\mat{A}$ singular, $\rg(\mat{A})<p$ y autovalores
        positivos con al menos uno nulo.
\end{itemize}
\end{mdframed}

\begin{observacion}
La clasificación positiva/semidefinida no es solo algebraica: determina si una
distancia o una varianza generalizada están bien definidas y no negativas.
\end{observacion}

\begin{ejemplo}
$\mat{A}=\begin{pmatrix}1&0\\0&1\end{pmatrix}$,
$Q=\mathbf{x}'\mat{A}\mathbf{x}=x_1^2+x_2^2>0$,
$|\mat{A}|=1$, $\rg(\mat{A})=2=p$,
$\lambda_1=\lambda_2=1>0$ $\Rightarrow$ \textbf{definida positiva}.
\end{ejemplo}

\begin{ejemplo}
$\mat{B}=\begin{pmatrix}1&-1\\-1&1\end{pmatrix}$,
$Q=x_1^2+x_2^2-2x_1x_2=(x_1-x_2)^2\ge0$,
$|\mat{B}|=0$, $\rg(\mat{B})=1<p$,
$\lambda_1=0$, $\lambda_2=2$ $\Rightarrow$ \textbf{semidefinida positiva}.
\end{ejemplo}

% ══════════════════════════════════════════════════════════════════════════
%  TEMA 2
% ══════════════════════════════════════════════════════════════════════════
\iniciotema{2}{Distribuciones Multivariantes}{Distribuciones Multivariantes}

% ══════════════════════════════════════════════════════════════════════════
\section{Introducción}
% ══════════════════════════════════════════════════════════════════════════

El \textbf{problema central del análisis de datos} es decidir si las
propiedades encontradas en una muestra pueden generalizarse a la población de
la que provienen. Para ello es necesario disponer de un \emph{modelo del
sistema generador de los datos}: suponer una distribución de probabilidad para
la variable aleatoria en la población.

En consecuencia, el primer paso del tema es definir con precisión distribución
conjunta, marginal y condicionada, y después estudiar cómo se traducen esas
estructuras en medidas de dependencia y en modelos paramétricos.

\begin{definicion}[Variable aleatoria multidimensional o vectorial]
Una variable aleatoria multidimensional es el resultado de observar
$p$ características en un elemento de la población. Puede ser:
\begin{itemize}
  \item \textbf{Discreta:} cuando lo es cada una de las $p$ variables
        escalares que la componen.
  \item \textbf{Continua:} cuando lo es cada una de las $p$ variables
        escalares que la componen.
  \item \textbf{Mixta:} cuando algunos componentes son discretos y otros
        continuos.
\end{itemize}
\end{definicion}

% ══════════════════════════════════════════════════════════════════════════
\section{Tipos de distribución}
% ══════════════════════════════════════════════════════════════════════════

\subsection{Distribución conjunta}

\begin{definicion}[Función de distribución conjunta]
La distribución de probabilidad que describe simultáneamente el comportamiento
del vector aleatorio $\mathbf{X}=(X_1,\dots,X_p)'$ es la \textbf{distribución
conjunta}. Dado $\mathbf{x}^0=(x_1^0,\dots,x_p^0)'$, su función de
distribución conjunta se define como:
\begin{mdframed}[style=resultado]
\[
  F_{\mathbf{X}}(\mathbf{x}^0) = P(\mathbf{X} \le \mathbf{x}^0)
  = P(X_1 \le x_1^0, \dots, X_p \le x_p^0).
\]
\end{mdframed}
$P(\mathbf{X}\le\mathbf{x}^0)$ representa la probabilidad acumulada sobre la
región ortante inferior determinada por $\mathbf{x}^0$. La función de
distribución conjunta es no decreciente en cada coordenada y contiene toda la
información probabilística del vector aleatorio.
\begin{variablesprimera}
  \varitem{\mathbf{X}}{vector aleatorio de dimensión $p$}
  \varitem{\mathbf{x}^0}{valor concreto del vector en el que se evalúa la distribución}
  \varitem{F_{\mathbf{X}}(\mathbf{x}^0)}{función de distribución conjunta acumulada}
  \varitem{P(\cdot)}{probabilidad del evento indicado}
\end{variablesprimera}
\end{definicion}

Habitualmente se trabaja con:
\begin{itemize}
  \item \textbf{Variable discreta:} función de probabilidad
        $p_{\mathbf{X}}(\mathbf{x}^0) = P(\mathbf{X}=\mathbf{x}^0)
        = P(X_1=x_1^0, \dots, X_p=x_p^0)$.
  \item \textbf{Variable continua:} el vector $\mathbf{X}$ es absolutamente
        continuo si existe una función de densidad $f_{\mathbf{X}}(\mathbf{x})$ tal que
        \[
          F_{\mathbf{X}}(\mathbf{x}^0)
          = \int_{(-\infty,\mathbf{x}^0]} f_{\mathbf{X}}(\mathbf{x})\,d\mathbf{x},
          \qquad d\mathbf{x} = dx_1\cdots dx_p.
        \]
\end{itemize}

La distribución conjunta es el objeto completo; las marginales y condicionadas
se obtienen de ella por integración y cociente, respectivamente.

\subsection{Distribuciones marginales}

\begin{definicion}[Distribución marginal]
Dada una variable aleatoria vectorial $p$-dimensional
$\mathbf{X}=(X_1,\dots,X_p)'$,
llamamos \textbf{distribución marginal} de cada componente a la distribución
univariante de dicho componente, sin tener en cuenta los otros. Para $p=2$:
\[
  f_{X_1}(x_1) = \int_{-\infty}^{+\infty} f_{X_1,X_2}(x_1,x_2)\,dx_2, \qquad
  f_{X_2}(x_2) = \int_{-\infty}^{+\infty} f_{X_1,X_2}(x_1,x_2)\,dx_1.
\]
\end{definicion}

Las marginales describen el comportamiento de cada componente por separado; la
conjunta añade, además, la estructura de dependencia entre componentes. Pasar de
conjunta a marginal simplifica descripción, pero elimina información de
interacción entre variables.

\subsection{Distribuciones condicionadas}

\begin{definicion}[Distribución condicionada]
Si $\mathbf{X}_1$ y $\mathbf{X}_2$ son subvectores de $\mathbf{X}$, la
\textbf{distribución condicionada} de $\mathbf{X}_1$ dado
$\mathbf{X}_2=\mathbf{x}_2^0$ se define como:
\[
  f_{\mathbf{X}_1 \mid \mathbf{X}_2}(\mathbf{x}_1 \mid \mathbf{x}_2^0)
  = \frac{f_{\mathbf{X}_1,\mathbf{X}_2}(\mathbf{x}_1, \mathbf{x}_2^0)}
         {f_{\mathbf{X}_2}(\mathbf{x}_2^0)},
  \qquad \text{si } f_{\mathbf{X}_2}(\mathbf{x}_2^0)\neq 0.
\]
\begin{variablesprimera}
  \varitem{\mathbf{X}_1}{subvector aleatorio de interés}
  \varitem{\mathbf{X}_2}{subvector aleatorio sobre el que se condiciona}
  \varitem{\mathbf{x}_2^0}{valor fijado del subvector $\mathbf{x}_2$}
  \varitem{f_{\mathbf{X}_1 \mid \mathbf{X}_2}(\mathbf{x}_1 \mid \mathbf{x}_2^0)}{densidad condicionada}
\end{variablesprimera}
\end{definicion}

La distribución condicionada cuantifica cómo cambia la ley de un bloque de
variables al fijar otro bloque: es la herramienta natural para inferencia
predictiva y para reglas de clasificación bayesianas.

% ══════════════════════════════════════════════════════════════════════════
\section{Independencia}
% ══════════════════════════════════════════════════════════════════════════

\begin{definicion}[Independencia]
Dos vectores aleatorios $\mathbf{X}_1$, $\mathbf{X}_2$ son
\textbf{independientes} si el conocimiento de uno no aporta información
respecto a los valores del otro:
\[
  f_{\mathbf{X}_2\mid\mathbf{X}_1}(\mathbf{x}_2\mid\mathbf{x}_1)
  = \frac{f_{\mathbf{X}_1,\mathbf{X}_2}(\mathbf{x}_1,\mathbf{x}_2)}
         {f_{\mathbf{X}_1}(\mathbf{x}_1)}
  = f_{\mathbf{X}_2}(\mathbf{x}_2),
\]
es decir, si
$f_{\mathbf{X}_1,\mathbf{X}_2}(\mathbf{x}_1,\mathbf{x}_2)
=f_{\mathbf{X}_1}(\mathbf{x}_1)f_{\mathbf{X}_2}(\mathbf{x}_2)$.
\end{definicion}

En general, las variables $\mathbf{X}_1, \dots, \mathbf{X}_p$ son
independientes si:
\[
  f_{\mathbf{X}_1,\dots,\mathbf{X}_p}(\mathbf{x}_1, \dots, \mathbf{x}_p)
  = \prod_{j=1}^p f_{\mathbf{X}_j}(\mathbf{x}_j).
\]
La independencia conjunta es una condición muy fuerte: cualquier subconjunto
de variables mantiene la independencia. Además, al aplicar funciones adecuadas
componente a componente, las variables transformadas siguen siendo independientes.

\begin{observacion}
Cuando las variables son independientes, no ganamos nada con su estudio
conjunto, y conviene estudiarlas individualmente.
\end{observacion}

\begin{observacion}
Independencia e incorrelación no son equivalentes en general. La equivalencia
solo se recupera bajo hipótesis adicionales (por ejemplo, normalidad
multivariante), que se explicitarán más adelante.
\end{observacion}

% ══════════════════════════════════════════════════════════════════════════
\section{Medidas características}
% ══════════════════════════════════════════════════════════════════════════

\subsection{Tendencia central: vector de medias}

\begin{definicion}[Vector de medias]
La \textbf{esperanza} o \textbf{vector de medias} $\boldsymbol{\mu}$ de una
variable multidimensional $\mathbf{X}$ es el vector cuyas componentes son las
esperanzas de las componentes:
\begin{mdframed}[style=resultado]
\[
  \boldsymbol{\mu} = E(\mathbf{X}), \qquad
  \bar{\mathbf{x}} = \frac{1}{n}\sum_{i=1}^n \mathbf{x}_i, \qquad
  \boldsymbol{\mu} = \int_{\R^p} \mathbf{x}\,f_{\mathbf{X}}(\mathbf{x})\,d\mathbf{x}.
\]
En notación matricial (con $\mat{X}\in\R^{n\times p}$), el vector de medias
muestral es $\bar{\mathbf{x}} = \dfrac{1}{n}\mat{X}'\mathbf{1}$.
\end{mdframed}
\begin{variablesprimera}
  \varitem{\boldsymbol{\mu}}{vector de medias poblacionales}
  \varitem{\mathbf{X}}{vector aleatorio multivariante}
  \varitem{\mathbf{x}_i}{observación $i$ de la muestra}
  \varitem{\bar{\mathbf{x}}}{vector de medias muestrales}
  \varitem{\mathbf{1}}{vector columna de unos de dimensión compatible}
\end{variablesprimera}
\end{definicion}

Propiedades del vector de medias:
\begin{itemize}
  \item Es el centro de gravedad euclídeo de la nube de puntos:
        haciendo nula la suma de desviaciones:
        $\displaystyle\sum_{i=1}^n (\mathbf{x}_i - \bar{\mathbf{x}})=\mathbf{0}$.
  \item Es un \textbf{operador lineal}:
        $E[\mat{A}\mathbf{X}+\mathbf{b}] = \mat{A}E[\mathbf{X}]+\mathbf{b}$.
\end{itemize}

En interpretación geométrica, $\boldsymbol{\mu}$ fija la localización de la
nube de puntos; la matriz de covarianzas fijará su forma, orientación y escala
en cada dirección.

\subsection{Variabilidad: matriz de varianzas y covarianzas}

\begin{definicion}[Matriz de varianzas y covarianzas]
En variables multivariantes, la variabilidad y la relación lineal entre
variables se presenta de forma compacta en la \textbf{matriz de varianzas y
covarianzas}:
\begin{mdframed}[style=resultado]
\[
  \mat{S}_x = E\bigl[(\mathbf{X}-\boldsymbol{\mu})(\mathbf{X}-\boldsymbol{\mu})'\bigr]
  = \frac{1}{n}\sum_{i=1}^n (\mathbf{x}_i-\bar{\mathbf{x}})(\mathbf{x}_i-\bar{\mathbf{x}})'.
\]
\end{mdframed}
\begin{variablesprimera}
  \varitem{\mat{S}_x}{matriz de varianzas y covarianzas}
  \varitem{S_i^2}{varianza de la variable $i$}
  \varitem{S_{ij}}{covarianza entre las variables $i$ y $j$}
\end{variablesprimera}
Explícitamente:
\[
  s_{jk} = \frac{1}{n}\sum_{i=1}^n (x_{ij}-\bar{x}_j)(x_{ik}-\bar{x}_k),
  \qquad j,k=1,\dots,p,
\]
\[
  \mat{S}_x =
  \begin{pmatrix}
    s_{11} & \cdots & s_{1p} \\
    \vdots & \ddots & \vdots \\
    s_{p1} & \cdots & s_{pp}
  \end{pmatrix},
  \qquad s_{jj}=s_j^2.
\]
\end{definicion}

\begin{observacion}
La diagonal de $\mat{S}_x$ mide dispersión individual y los términos
extradiagonales cuantifican dependencia lineal entre variables.
\end{observacion}

\subsubsection{Matriz de datos centrados}

Se define la \textbf{matriz de datos centrados} (o puntuaciones diferenciales)
$\tilde{\mathbf{X}}$ como la matriz resultado de restar a cada dato su media:
\[
  \tilde{\mathbf{X}} = \mat{X} - \mathbf{1}\bar{\mathbf{x}}'
  = \mat{X} - \frac{1}{n}\mathbf{1}\mathbf{1}'\mat{X}
  = \mat{P}\mat{X},
\]
y la \textbf{matriz de proyección} es:
\[
  \mat{P} = \mat{I} - \frac{1}{n}\mathbf{1}\cdot\mathbf{1}',
\]
que es simétrica ($\mat{P}=\mat{P}'$) e \textbf{idempotente} ($\mat{P}^2=\mat{P}$).
\begin{variablesprimera}
  \varitem{\tilde{\mathbf{X}}}{matriz de datos centrados}
  \varitem{\mat{P}}{matriz de proyección de centrado}
  \varitem{\mat{I}}{matriz identidad de orden $n$}
\end{variablesprimera}

Por tanto:
\[
  \mat{S} = \frac{1}{n}\tilde{\mathbf{X}}'\tilde{\mathbf{X}}
  = \frac{1}{n}\mat{X}'\mat{P}\mat{X}.
\]

La \textbf{matriz de varianzas corregida} (estimación insesgada) es:
\[
  \hat{\mat{S}} = \frac{1}{n-1}\tilde{\mathbf{X}}'\cdot\tilde{\mathbf{X}}.
\]

\subsubsection{S es semidefinida positiva}

\begin{proposicion}
La matriz $\mat{S}$ es semidefinida positiva, es decir,
$\mathbf{y}'\mat{S}\mathbf{y}\ge0$ para todo $\mathbf{y}\in\R^p$.
\end{proposicion}

\begin{proof}
Sea $\mathbf{w}\in\R^p$ y definamos la combinación lineal centrada
$v_i=\mathbf{w}'(\mathbf{x}_i-\bar{\mathbf{x}})$. Su media es nula y su
varianza muestral:
\[
  \Var(v) = \frac{1}{n}\sum_{i=1}^n v_i^2
  = \frac{1}{n}\sum_{i=1}^n [\mathbf{w}'(\mathbf{x}_i-\bar{\mathbf{x}})]
    [(\mathbf{x}_i-\bar{\mathbf{x}})'\mathbf{w}] = \mathbf{w}'\mat{S}\mathbf{w} \ge 0.
\]
Como toda varianza es no negativa y $\mathbf{w}$ es arbitrario,
$\mat{S}$ es semidefinida positiva.
\end{proof}

\begin{observacion}
Este resultado no es accesorio: garantiza varianzas no negativas para toda
combinación lineal y asegura autovalores reales no negativos, base matemática
de ACP, Mahalanobis y discriminante gaussiano.
\end{observacion}

Consecuencias:
\begin{itemize}
  \item Si $\mat{S}\mathbf{w}_i = \lambda_i\mathbf{w}_i$: $\lambda_i\ge0$,
        todos los menores principales son no negativos y $|\mat{S}|\ge0$.
  \item Si $\mat{S}$ tiene rango $h<p$, existen $r=p-h$ combinaciones lineales
        entre las variables, es decir $r$ variables redundantes que pueden
        eliminarse. Se pueden representar las observaciones con $h=p-r$ variables.
  \item Cualquier vector del subespacio definido por $(\mathbf{w}_1,\dots,\mathbf{w}_r)$
        puede expresarse como combinación lineal y verifica
        $\mat{S}(a_1\mathbf{w}_1+\cdots+a_r\mathbf{w}_r)=\mathbf{0}$.
        Los autovectores de $\mat{S}$ asociados a valores propios nulos
        constituyen una base ortonormal en dicho espacio.
\end{itemize}

\subsubsection{Variabilidad total}

\begin{definicion}[Variabilidad total]
Una forma de resumir la variabilidad de un conjunto de variables es mediante la
traza de su matriz de varianzas y covarianzas:
\[
  T = \tr(\mat{S}) = \sum_{j=1}^p s_j^2.
\]
\end{definicion}

\begin{observacion}
Una dependencia alta entre variables implica una variabilidad conjunta pequeña.
Esta medida prescinde de las relaciones de dependencia existentes.
\end{observacion}

\subsubsection{Varianza generalizada}

\begin{definicion}[Varianza generalizada]
\[
  VG = |\mat{S}|.
\]
Mide mejor la variabilidad global. Su raíz cuadrada se denomina
\textbf{desviación típica generalizada}, y tiene las siguientes propiedades:
\begin{itemize}
  \item Está bien definida (el determinante de la matriz de varianzas y
        covarianzas es siempre no negativo).
  \item Es una medida del área ($p=2$), volumen ($p=3$) o hipervolumen
        ($p>3$) ocupado por el conjunto de datos.
\end{itemize}
\end{definicion}

Para $p=2$:
\[
  \mat{S} = \begin{pmatrix}s_x^2 & rs_xs_y \\ rs_xs_y & s_y^2\end{pmatrix},
  \qquad |\mat{S}|^{1/2} = s_xs_y\sqrt{1-r^2}.
\]
Si las variables están relacionadas linealmente ($r\neq0$), la mayoría de los
puntos tenderán a situarse en una franja alrededor de la recta de regresión,
reduciéndose el área tanto mayor cuanto mayor sea $r^2$.

\subsubsection{Distancias}

\begin{definicion}[Distancia o métrica]
Dados dos puntos $\mathbf{x}_i, \mathbf{x}_j\in\R^p$, una \textbf{distancia} es
una función $d(\mathbf{x}_i,\mathbf{x}_j)$ que satisface:
\begin{enumerate}
  \item $d(\mathbf{x}_i,\mathbf{x}_j)\ge0$ (no negatividad).
  \item $d(\mathbf{x}_i,\mathbf{x}_i)=0$ (identidad).
  \item $d(\mathbf{x}_i,\mathbf{x}_j)=d(\mathbf{x}_j,\mathbf{x}_i)$ (simetría).
  \item $d(\mathbf{x}_i,\mathbf{x}_j)\le d(\mathbf{x}_i,\mathbf{x}_\ell)+d(\mathbf{x}_\ell,\mathbf{x}_j)$ (desigualdad triangular).
\end{enumerate}
\end{definicion}

La distancia euclídea es la referencia geométrica básica, pero depende de
escala y no incorpora correlaciones. Por ello, en multivariante se emplean
también distancias ponderadas y, en particular, la de Mahalanobis.

\begin{definicion}[Distancia de Mahalanobis]
La distancia de Mahalanobis entre un punto $\mathbf{x}_i$ y el vector de
medias $\bar{\mathbf{x}}$ se define como:
\begin{mdframed}[style=resultado]
\[
  d_M = \sqrt{(\mathbf{x}_i-\bar{\mathbf{x}})'\mat{S}^{-1}(\mathbf{x}_i-\bar{\mathbf{x}})}.
\]
\end{mdframed}
Tiene en cuenta variabilidad y correlación, y es invariante frente a cambios
lineales de coordenadas no singulares.
\end{definicion}

\begin{observacion}
Interpretación: la distancia euclídea mide separación en coordenadas
originales, mientras que Mahalanobis mide separación en la geometría inducida
por $\mat{S}$; por eso penaliza menos direcciones de alta varianza y más las de
baja varianza.
\end{observacion}

Para $p=2$, siendo $s_{1,2}=rs_1s_2$:
\[
  \mat{S}^{-1} = \frac{1}{1-r^2}
  \begin{pmatrix}s_1^{-2}&-rs_1^{-1}s_2^{-1}\\-rs_1^{-1}s_2^{-1}&s_2^{-2}\end{pmatrix},
\]
\[
  d_M^2 = \frac{1}{1-r^2}\left[
    \frac{\Delta x^2}{s_1^2} + \frac{\Delta y^2}{s_2^2}
    - 2r\frac{\Delta x\,\Delta y}{s_1s_2}
  \right].
\]
\[
  \text{donde } \Delta x=x_{i1}-\bar{x}_1,\quad \Delta y=x_{i2}-\bar{x}_2.
\]

\begin{itemize}
  \item Si $r=0$: se reduce a la distancia euclídea estandarizada.
  \item Si $r$ es baja: se aproxima a la euclídea y cada variable contribuye
        de forma independiente.
  \item Si $r$ es alta: es menor que la euclídea y $\mat{S}$ reduce la
        contribución de las variables correlacionadas.
  \item Si $r=1$: $\mat{S}$ es singular, una variable es redundante y la
        distancia es indefinida.
\end{itemize}

\subsubsection{Distancia promedio}

Podemos medir la variabilidad promediando las distancias entre los puntos y la
media. Con distancias euclídeas:
\[
  V_m = \frac{1}{n}\sum_{i=1}^n (\mathbf{x}_i-\bar{\mathbf{x}})'(\mathbf{x}_i-\bar{\mathbf{x}})
  = \sum_{i=1}^n \tr\!\left[\frac{1}{n}(\mathbf{x}_i-\bar{\mathbf{x}})(\mathbf{x}_i-\bar{\mathbf{x}})'\right]
  = \tr(\mat{S}).
\]
El promedio de distancias es la variabilidad total. Si se promedia también por
la dimensión del vector:
\[
  V_{m,p} = \frac{1}{np}\sum_{i=1}^n (\mathbf{x}_i-\bar{\mathbf{x}})'(\mathbf{x}_i-\bar{\mathbf{x}}) = \bar{s}^2.
\]

% ══════════════════════════════════════════════════════════════════════════
\section{Medidas de dependencia lineal}
% ══════════════════════════════════════════════════════════════════════════

El objetivo fundamental de la descripción de datos multivariantes es
comprender la estructura de dependencias entre las variables. Estas
dependencias pueden estudiarse:
\begin{itemize}
  \item Entre pares de variables.
  \item Entre una variable y el resto.
  \item Entre pares de variables eliminando el efecto de las demás.
  \item Entre el conjunto de todas las variables.
\end{itemize}
En esta sección se pasa de dependencia bivariante a dependencia condicional y,
finalmente, a medidas globales del sistema completo.

\subsection{Dependencia entre pares: coeficiente de correlación}

\begin{definicion}[Coeficiente de correlación lineal]
\[
  r_{jk} = \frac{s_{jk}}{s_j s_k}, \qquad 0\le|r_{jk}|\le1.
\]
Si existe relación lineal exacta ($x_{ij}=a+bx_{ik}$): $|r_{jk}|=1$. Es
invariante ante transformaciones lineales de las variables.
\end{definicion}

\begin{definicion}[Matriz de correlaciones]
La \textbf{matriz de correlaciones} del vector aleatorio $\mathbf{x}$ con
matriz de covarianzas $\mat{S}_x$ es:
\begin{mdframed}[style=resultado]
\[
  \mat{R}_x = \mat{D}^{-\frac{1}{2}}\mat{S}_x\mat{D}^{-\frac{1}{2}},
  \qquad \mat{D} = \diag(\sigma_1^2,\dots,\sigma_p^2).
\]
\end{mdframed}
\[
  \mat{R} = \begin{pmatrix}1&\cdots&r_{1p}\\\vdots&\ddots&\vdots\\r_{p1}&\cdots&1\end{pmatrix}.
\]
Es una matriz cuadrada y simétrica con unos en la diagonal y coeficientes de
correlación fuera. Es semidefinida positiva.
\end{definicion}

\begin{observacion}
La fórmula anterior se obtiene al estandarizar cada variable por su desviación
típica: si $\mathbf{Z}=(\mathbf{X}-\boldsymbol{\mu})\mat{D}^{-1/2}$, entonces
$\Cov(\mathbf{Z})=\mat{R}_x$.
\end{observacion}

\begin{proposicion}
La matriz de correlaciones $\mat{R}_x$ es simétrica y semidefinida positiva.
\end{proposicion}

\begin{proof}
La simetría se hereda de $\mat{S}_x$. Además, para cualquier
$\mathbf{u}\in\R^p$:
\[
  \mathbf{u}'\mat{R}_x\mathbf{u}
  = (\mat{D}^{-1/2}\mathbf{u})'\mat{S}_x(\mat{D}^{-1/2}\mathbf{u})\ge 0,
\]
porque $\mat{S}_x$ es semidefinida positiva. Por tanto, $\mat{R}_x$ también
es semidefinida positiva.
\end{proof}

\subsection{Dependencia entre una variable y el resto: regresión múltiple}

Se estudia mediante \textbf{regresión múltiple}. El predictor lineal de la
variable $y$ es:
\[
  \hat{y}_i = \beta_0 + \beta_1 x_{i1} + \beta_2 x_{i2} + \cdots + \beta_k x_{ik} + e_i,
\]
donde $e_i = y_i - \hat{y}_i$ son los errores de predicción. El criterio de
mínimos cuadrados asigna a los $\beta_j$ el valor que minimiza la suma de
errores al cuadrado. Los residuos deben estar incorrelados con las variables
explicativas (\emph{covarianza cero}): el vector de residuos es ortogonal al
espacio generado por las variables explicativas.

\subsection{Correlación parcial: eliminando el efecto de las demás}

\begin{definicion}[Coeficiente de correlación parcial]
El coeficiente de correlación parcial entre $(x_1,x_2)$ dadas
$(x_3,\dots,x_p)$ se denota $r_{12.3\dots p}$: es el coeficiente de
correlación entre $x_1$ y $x_2$ que están libres de los efectos de
$(x_3,\dots,x_p)$.

Se calcula:
\begin{enumerate}[label=\alph*)]
  \item Se obtiene la parte de cada variable no explicada por el grupo de
        variables controladas: \emph{residuo de la regresión sobre
        $(x_3,\dots,x_p)$}.
  \item Se calcula el coeficiente de correlación simple entre estos dos
        residuos.
\end{enumerate}
Puede demostrarse que estandarizando los elementos de $\mat{S}^{-1}$, llamando
$s^{ij}$ a sus elementos:
\begin{mdframed}[style=resultado]
\[
  r_{jk.\,\text{resto}} = -\frac{s^{jk}}{\sqrt{s^{jj}s^{kk}}}.
\]
\end{mdframed}
\end{definicion}

\begin{observacion}
La expresión anterior muestra que la correlación parcial es nula exactamente
cuando el término fuera de la diagonal correspondiente en la matriz de
precisión $\mat{S}^{-1}$ es cero; este hecho conecta la dependencia condicional
lineal con la estructura de la inversa de la covarianza.
\end{observacion}

La \textbf{matriz de correlaciones parciales} $\mat{P}$ contiene los
coeficientes de correlación parcial entre pares de variables eliminando el
efecto de las restantes. Por ejemplo, para cuatro variables:
\[
  \mat{P}_4 = \begin{pmatrix}
    1 & r_{12.34} & r_{13.24} & r_{14.23} \\
    r_{21.34} & 1 & r_{23.14} & r_{24.13} \\
    r_{31.24} & r_{32.14} & 1 & r_{34.12} \\
    r_{41.23} & r_{42.13} & r_{43.12} & 1
  \end{pmatrix}.
\]

\subsection{Dependencia global: determinante de \texorpdfstring{$\mat{R}$}{R}}

\begin{definicion}[Medida global de dependencia]
El \textbf{determinante de la matriz de correlación} mide lo alejado que está
el conjunto de variables de la situación de dependencia lineal perfecta:
\[
  0\le|\mat{R}|\le1.
\]
\begin{itemize}
  \item Si las variables están todas incorreladas: $\mat{R}$ es diagonal con
        unos y $|\mat{R}|=1$.
  \item Si una variable es combinación lineal del resto: $\mat{S}$, $\mat{R}$
        son singulares y $|\mat{R}|=0$.
\end{itemize}
\end{definicion}

\begin{definicion}[Dependencia $D_p$ (Peña y Rodríguez, 2000)]
Una medida global de las correlaciones lineales es:
\[
  D_p = 1 - |\mat{R}_p|^{1/(p-1)}.
\]
Para $p=2$: $|\mat{R}_p|=1-r_{12}^2$, luego $D_p = r_{12}^2$, que coincide
con el coeficiente de determinación.
\end{definicion}

\subsection{Matriz de precisión}

\begin{definicion}[Matriz de precisión]
La \textbf{matriz de precisión} es la inversa de la matriz de varianzas y
covarianzas: $\mat{S}^{-1}$. Contiene la información sobre la relación
multivariante entre cada variable y el resto.
\end{definicion}

Sus elementos tienen la siguiente interpretación:
\begin{itemize}
  \item \textbf{Fuera de la diagonal:}
        $s^{ij} = -\hat{\beta}_{ij}/s_r^2(i)$,
        donde $\hat{\beta}_{ij}$ es el coeficiente de regresión de la variable
        $j$ para explicar la variable $i$, y $s_r^2(i)$ es la varianza
        residual de esa regresión.
  \item \textbf{En la diagonal:} $s^{ii} = 1/s_r^2(i)$ (inversas de las
        varianzas residuales de cada variable en su regresión con el resto).
  \item Si se estandarizan los elementos para que tenga unos en la diagonal,
        los elementos fuera de la diagonal son los coeficientes de correlación
        parcial:
        $r_{ij,R} = -s^{ij}/\sqrt{s^{ii}s^{jj}}$,
        donde $R$ denota el conjunto de $p-2$ variables $x_k$ con $k\neq i,j$.
\end{itemize}

\begin{observacion}
En términos de modelización probabilística, la matriz de precisión resume la
estructura de dependencia condicional lineal; por ello es central en modelos
gaussianos, redes gaussianas y reglas discriminantes lineales.
\end{observacion}

\begin{ejemplo}
\[
  \mat{S}^{-1} = \begin{pmatrix}
    52{,}0942 & -47{,}9058 & 52{,}8796 \\
    -47{,}9058 & 52{,}0942 & -47{,}1204 \\
    52{,}8796 & -47{,}1204 & 60{,}2094
  \end{pmatrix}.
\]
La primera fila puede escribirse como $52{,}0942\times(1,-0{,}9196,1{,}0151)$:
los coeficientes de regresión de $x_2$ y $x_3$ para explicar $x_1$ son
$0{,}9196$ y $-1{,}0151$ respectivamente. La varianza residual de la regresión
de $x_1$ sobre las otras dos es $1/52{,}0942 = 0{,}0192$.
\end{ejemplo}

% ══════════════════════════════════════════════════════════════════════════
\section{Transformaciones de vectores aleatorios}
% ══════════════════════════════════════════════════════════════════════════

\subsection{Generalidades: función de densidad y cambio de variable}

\begin{definicion}[Función de densidad]
La función de densidad describe la densidad de probabilidad en cualquier
intervalo:
\[
  f(x_0)=\lim_{\Delta x\to 0^+}
  \frac{P(x_0 < X \le x_0+\Delta x)}{\Delta x}.
\]
Al hacer las clases del histograma cada vez más pequeñas, éste tenderá a una
curva $f(x)$ que describe el comportamiento de la variable.
\end{definicion}

\paragraph{Propiedades de $f(x)$:}
\begin{enumerate}
  \item $f(x)\ge0\quad\forall x\in\R$.
  \item $P(a<x\le b) = \displaystyle\int_a^b f(x)\,dx\quad\forall a,b\in\R$.
  \item $P(X\le x) = \displaystyle\int_{-\infty}^x f(x)\,dx$.
  \item $\displaystyle\int_{-\infty}^{+\infty}f(x)\,dx = 1$.
\end{enumerate}

Como la función de densidad tiene dimensiones, si cambian las unidades de
medida, la función debe modificarse también.

\paragraph{Cambio de variable general.}
Sea $\mathbf{x}\in\R^p$ con función de densidad $f_x(\mathbf{x})$ y sea otro
vector aleatorio $\mathbf{y}\in\R^p$ definido por la transformación:
\[
  y_i = g_i(x_1,\dots,x_p), \quad i=1,\dots,p.
\]
Suponiendo que existen las funciones inversas $x_i = h_i(y_1,\dots,y_p)$ y
que todas son diferenciables, la transformación es biyectiva en el soporte y
el jacobiano no se anula, la función de densidad de $\mathbf{y}$ es:
\begin{mdframed}[style=resultado]
\[
  f_y(\mathbf{y}) = f_x(\mathbf{x})\,\left|\frac{d\mathbf{x}}{d\mathbf{y}}\right|,
\]
\end{mdframed}
donde $\left|\dfrac{d\mathbf{x}}{d\mathbf{y}}\right|$ es el \textbf{jacobiano}
de la transformación:
\[
  \left|\frac{d\mathbf{x}}{d\mathbf{y}}\right|
  = \begin{vmatrix}
      \partial x_1/\partial y_1 & \cdots & \partial x_1/\partial y_p \\
      \vdots & & \vdots \\
      \partial x_p/\partial y_1 & \cdots & \partial x_p/\partial y_p
    \end{vmatrix}.
\]

\paragraph{Caso de transformaciones lineales.}
Para $\mathbf{y}=\mat{A}\mathbf{x}$ con $\mat{A}$ cuadrada no singular:
$\mathbf{x}=\mat{A}^{-1}\mathbf{y}$, luego $\partial\mathbf{x}/\partial\mathbf{y}=\mat{A}^{-1}$
y el jacobiano es $|\mat{A}^{-1}|=|\mat{A}|^{-1}$. Por tanto:
\begin{mdframed}[style=resultado]
\[
  f_y(\mathbf{y}) = f_x\!\left(\mat{A}^{-1}\mathbf{y}\right)|\mat{A}|^{-1}.
\]
\end{mdframed}
Para obtener la función de densidad de $\mathbf{y}$, se sustituye el argumento
en $f_x$ por $\mat{A}^{-1}$ y se divide por $|\mat{A}|$.

\subsection{Esperanzas de transformaciones lineales}

Sea $\mathbf{X}$ un vector aleatorio de dimensión $p$ y $\mathbf{Y}=\mat{A}\mathbf{X}$
un nuevo vector de dimensión $m\le p$ ($\mat{A}$ rectangular de dimensiones
$m\times p$). Llamando $\boldsymbol{\mu}_x, \boldsymbol{\mu}_y$ a los vectores
de medias y $\mat{S}_x, \mat{S}_y$ a las matrices de covarianzas:

\begin{mdframed}[style=resultado]
\[
  \boldsymbol{\mu}_y = \mat{A}\boldsymbol{\mu}_x, \qquad
  \mat{S}_y = \mat{A}\mat{S}_x\mat{A}'.
\]
\end{mdframed}

\begin{proof}
Por linealidad de la esperanza:
\[
  \boldsymbol{\mu}_y=E(\mathbf{Y})=E(\mat{A}\mathbf{X})=\mat{A}E(\mathbf{X})
  =\mat{A}\boldsymbol{\mu}_x.
\]
Además,
\[
  \mat{S}_y
  = \Cov(\mathbf{Y})
  = E\!\left[(\mathbf{Y}-\boldsymbol{\mu}_y)(\mathbf{Y}-\boldsymbol{\mu}_y)'\right]
  = E\!\left[\mat{A}(\mathbf{X}-\boldsymbol{\mu}_x)(\mathbf{X}-\boldsymbol{\mu}_x)'\mat{A}'\right].
\]
Como $\mat{A}$ es constante respecto a la esperanza:
\[
  \mat{S}_y
  = \mat{A}\,E\!\left[(\mathbf{X}-\boldsymbol{\mu}_x)(\mathbf{X}-\boldsymbol{\mu}_x)'\right]\mat{A}'
  = \mat{A}\mat{S}_x\mat{A}'.
\]
\end{proof}

\paragraph{Esperanzas condicionadas.}
\[
  E(\mathbf{X}_1\mid\mathbf{X}_2) = \int \mathbf{x}_1\,f(\mathbf{x}_1\mid\mathbf{x}_2)\,d\mathbf{x}_1.
\]
En general será función del valor de $\mathbf{X}_2$. Si $\mathbf{X}_2$ es
fija, la esperanza condicionada es una constante; si es variable aleatoria,
será una variable aleatoria.

% ══════════════════════════════════════════════════════════════════════════
\section{Modelos de distribuciones multivariantes}
% ══════════════════════════════════════════════════════════════════════════

\subsection{Distribución multinomial}

\begin{definicion}[Distribución multinomial]
Se observan elementos que se pueden clasificar en $G$ clases. Si el proceso
generador es estable se definen $G$ variables aleatorias:
\[
  x_j = \begin{cases}1 & \text{si la observación pertenece al grupo }j\\
                     0 & \text{en otro caso}\end{cases}, \quad j=1,\dots,G.
\]
El vector de probabilidades de pertenencia a las clases es
$\mathbf{p}=(p_1,\dots,p_G)'$ con $\sum p_j=1$.

El resultado de una observación es un valor del vector de $G$-variables
$\mathbf{x}=(x_1,\dots,x_G)'$ de la forma $\mathbf{x}=(0,\dots,1,\dots,0)$
(el valor 1 corresponde a la clase observada).
\end{definicion}

Propiedades:
\begin{itemize}
  \item $E[\mathbf{x}] = \mathbf{p}$.
  \item $\mat{S}_x = \operatorname{diag}(\mathbf{p}) - \mathbf{p}\mathbf{p}'$:
        las varianzas son $p_j(1-p_j)$ y las covarianzas $-p_jp_k$.
\end{itemize}

Si se realizan $n$ observaciones independientes y $n_j$ es el número de
observaciones en la clase $j$, la \textbf{distribución multinomial} de
$\mathbf{n}=(n_1,\dots,n_G)'$ es:
\[
  P(\mathbf{n}) = \frac{n!}{n_1!\cdots n_G!}\,p_1^{n_1}\cdots p_G^{n_G},
  \qquad \sum n_j=n.
\]

\subsection{Distribución de Dirichlet}

\begin{definicion}[Distribución de Dirichlet]
La \textbf{distribución de Dirichlet} es la generalización multivariante de
la distribución Beta. Se define sobre el símplex estándar
$\mathcal{S}_G = \{(x_1,\dots,x_G): x_i\ge0, \sum x_i=1\}$ con parámetros
$\boldsymbol{\alpha}=(\alpha_1,\dots,\alpha_G)'$, $\alpha_i>0$.

Su función de densidad es:
\begin{mdframed}[style=resultado]
\[
  f(x_1,\dots,x_{G-1}) = \frac{\Gamma(\alpha_0)}{\prod_{i=1}^G\Gamma(\alpha_i)}
  \prod_{i=1}^G x_i^{\alpha_i-1},
  \qquad \alpha_0 = \sum_{i=1}^G\alpha_i.
\]
\end{mdframed}
\end{definicion}

Propiedades:
\begin{itemize}
  \item $E[x_i] = \dfrac{\alpha_i}{\alpha_0}$.
  \item $\operatorname{Var}[x_i] = \dfrac{\alpha_i(\alpha_0-\alpha_i)}{\alpha_0^2(\alpha_0+1)}$.
  \item $\operatorname{Cov}[x_i,x_j] = -\dfrac{\alpha_i\alpha_j}{\alpha_0^2(\alpha_0+1)}$
        para $i\neq j$.
  \item Cada marginal $x_i\sim\operatorname{Beta}(\alpha_i,\alpha_0-\alpha_i)$.
  \item Es la distribución conjugada de la multinomial: si
        $\mathbf{p}\sim\operatorname{Dir}(\boldsymbol{\alpha})$ y se observan
        $n_j$ eventos de clase $j$, la distribución a posteriori es
        $\operatorname{Dir}(\alpha_1+n_1,\dots,\alpha_G+n_G)$.
\end{itemize}

\subsection{Distribución normal multivariante}

La distribución más importante en el análisis multivariante es la
\textbf{distribución normal multivariante}, base teórica de la mayoría de
las técnicas clásicas. Su centralidad se debe a estabilidad bajo
transformaciones lineales, tractabilidad analítica y relación directa con
distancias cuadráticas y reglas de decisión.

\begin{definicion}[Normal multivariante]
Un vector aleatorio $\mathbf{X}$ de dimensión $p$ sigue una distribución
normal multivariante $\mathbf{X}\sim\mathcal{N}_p(\boldsymbol{\mu},\boldsymbol{\Sigma})$
si su función de densidad es:
\begin{mdframed}[style=resultado]
\[
  f(\mathbf{x}) = \frac{1}{(2\pi)^{p/2}|\boldsymbol{\Sigma}|^{1/2}}
  \exp\!\left\{-\tfrac{1}{2}(\mathbf{x}-\boldsymbol{\mu})'\boldsymbol{\Sigma}^{-1}(\mathbf{x}-\boldsymbol{\mu})\right\},
  \quad \mathbf{x}\in\R^p,
\]
\end{mdframed}
\begin{variablesprimera}
  \varitem{\boldsymbol{\Sigma}}{matriz de covarianzas común, simétrica y definida positiva}
  \varitem{|\boldsymbol{\Sigma}|}{determinante de la matriz de covarianzas}
  \varitem{\boldsymbol{\Sigma}^{-1}}{matriz inversa de $\boldsymbol{\Sigma}$}
\end{variablesprimera}
\end{definicion}

\paragraph{Propiedades fundamentales.}

\begin{itemize}
  \item \textbf{Marginales:} toda marginal de una normal multivariante es
        normal univariante: $X_i\sim\mathcal{N}(\mu_i,\sigma_i^2)$.

  \item \textbf{Combinaciones lineales:} si $\mathbf{X}\sim\mathcal{N}_p(\boldsymbol{\mu},\boldsymbol{\Sigma})$
        y $\mathbf{Y}=\mat{A}\mathbf{X}+\mathbf{b}$, entonces
        $\mathbf{Y}\sim\mathcal{N}_m(\mat{A}\boldsymbol{\mu}+\mathbf{b},\mat{A}\boldsymbol{\Sigma}\mat{A}')$.

  \item \textbf{Independencia $\Leftrightarrow$ incorrelación:} en la normal
        multivariante, dos variables son independientes si y solo si su
        covarianza es nula. Equivalentemente, la matriz de covarianzas es
        diagonal si y solo si las componentes son independientes.

  \item \textbf{Distribución condicionada:} si particionamos
        $\mathbf{X}=(\mathbf{X}_1',\mathbf{X}_2')'$ con
        $\boldsymbol{\mu}=(\boldsymbol{\mu}_1',\boldsymbol{\mu}_2')'$ y
        $\boldsymbol{\Sigma}=\begin{pmatrix}\boldsymbol{\Sigma}_{11}&\boldsymbol{\Sigma}_{12}\\\boldsymbol{\Sigma}_{21}&\boldsymbol{\Sigma}_{22}\end{pmatrix}$,
        entonces:
        \[
          \mathbf{X}_1|\mathbf{X}_2=\mathbf{x}_2 \sim \mathcal{N}\!\left(
          \boldsymbol{\mu}_1 + \boldsymbol{\Sigma}_{12}\boldsymbol{\Sigma}_{22}^{-1}(\mathbf{x}_2-\boldsymbol{\mu}_2),\;
          \boldsymbol{\Sigma}_{11} - \boldsymbol{\Sigma}_{12}\boldsymbol{\Sigma}_{22}^{-1}\boldsymbol{\Sigma}_{21}
          \right).
        \]

  \item \textbf{Distancia de Mahalanobis:} la forma cuadrática
        $(\mathbf{X}-\boldsymbol{\mu})'\boldsymbol{\Sigma}^{-1}(\mathbf{X}-\boldsymbol{\mu})\sim\chi^2(p)$
        (distribución chi-cuadrado con $p$ grados de libertad). Las
        superficies de igual densidad son elipsoides centrados en
        $\boldsymbol{\mu}$.

  \item \textbf{Estimación:} los estimadores de máxima verosimilitud son
        $\hat{\boldsymbol{\mu}}=\bar{\mathbf{x}}$ y
        $\hat{\boldsymbol{\Sigma}}=\frac{1}{n}\sum(\mathbf{x}_i-\bar{\mathbf{x}})(\mathbf{x}_i-\bar{\mathbf{x}})'=\mat{S}$.
        El estimador insesgado de $\boldsymbol{\Sigma}$ es $\hat{\mat{S}}=\frac{n}{n-1}\mat{S}$.
\end{itemize}

% ══════════════════════════════════════════════════════════════════════════
%  TEMA 3
% ══════════════════════════════════════════════════════════════════════════
\iniciotema{3}{Análisis de Componentes Principales}{Análisis de Componentes Principales}

% ══════════════════════════════════════════════════════════════════════════
\section{Introducción}
% ══════════════════════════════════════════════════════════════════════════

El \textbf{Análisis de Componentes Principales (ACP)} es un tipo de análisis
factorial —junto con el Análisis Factorial de Correlaciones y el Análisis de
Correspondencias— cuyo objetivo es analizar la estructura de las correlaciones
entre un gran número de variables mediante la definición de una serie de
\emph{dimensiones subyacentes o factores}.

Formalmente, ACP construye una base ortogonal adaptada a la nube de datos:
ordena direcciones por varianza explicada y permite representar el problema en
dimensión reducida con pérdida controlada.

\begin{itemize}
  \item Es una \textbf{técnica de interdependencia}: se consideran todas las
        variables simultáneamente, cada una relacionada con todas las demás.
  \item Usa el \textbf{valor teórico} (combinación lineal de las variables).
  \item Los valores teóricos (factores o componentes) tratan de maximizar la
        explicación de la serie de variables, \emph{no} de predecir variables
        dependientes.
  \item Permite la \textbf{síntesis} de un gran conjunto de datos: la
        información contenida se representa mediante un conjunto de variables
        (componentes), función lineal de las originales, de número menor.
  \item Los componentes suelen interpretarse como \textbf{factores
        responsables de la variabilidad observada}.
\end{itemize}

La idea geométrica central es proyectar los datos sobre subespacios de menor
dimensión maximizando la varianza retenida (equivalentemente, minimizando error
cuadrático de reconstrucción). Es decir, ACP resuelve un problema simultáneo de
síntesis estadística y aproximación geométrica.

\begin{nota}
El origen de la técnica se debe a \textbf{Karl Pearson (1901)}, que la
introdujo como un medio de ajustar hiperplanos según el criterio de mínimos
cuadrados ortogonales. Fue establecida formalmente por \textbf{Harold
Hotelling (1933)} y desarrollada gracias a los avances informáticos del
siglo XX y XXI.
\end{nota}

Dado un conjunto de $n$ observaciones de $p$ variables, se trata de representar
adecuadamente esta información con un número menor de variables construidas
como combinaciones lineales de las originales:

\begin{mdframed}[style=resultado]
Con variables de alta dependencia es frecuente que un pequeño número de nuevas
variables ($<$20\,\% de las originales) expliquen la mayor parte ($>$80\,\%)
de la variabilidad original.
\end{mdframed}

La técnica permite:
\begin{itemize}
  \item Representar óptimamente observaciones de un espacio $p$-dimensional
        en un espacio de dimensión más reducida (primer paso para identificar
        variables ``latentes'' o no observadas que generan la variabilidad).
  \item Transformar las variables originales, en general correladas, en nuevas
        variables incorreladas, facilitando la interpretación de los datos.
\end{itemize}

\paragraph{Cuestiones prácticas.}
\begin{itemize}
  \item No deben usarse muestras inferiores a 50 observaciones (idealmente
        100 o más).
  \item Como regla general, la muestra mínima debería ser \textbf{cinco veces
        mayor} que el número de variables analizadas; el tamaño aceptable es
        un ratio de \textbf{diez a uno}.
  \item Si se incluyen indiscriminadamente grandes cantidades de variables
        esperando que el análisis factorial lo solucione, la posibilidad de
        obtener malos resultados es alta.
\end{itemize}

% ══════════════════════════════════════════════════════════════════════════
\section{Planteamiento del problema}
% ══════════════════════════════════════════════════════════════════════════

Partimos de variables $\mathbf{x}_1,\mathbf{x}_2,\dots,\mathbf{x}_p$ medidas
sobre $n$ unidades observacionales $\Rightarrow$ matriz $\mathbf{X}$ de orden
$n\times p$. Queremos unas nuevas variables latentes o \textbf{componentes}
$\mathbf{y}_1,\mathbf{y}_2,\dots,\mathbf{y}_p$ que:
\begin{variablesprimera}
  \varitem{\mathbf{y}_j}{componente principal $j$}
\end{variablesprimera}

\begin{enumerate}[label=\alph*)]
  \item No establezcan ningún requisito sobre la distribución de
        $\mathbf{x}_1,\mathbf{x}_2,\dots,\mathbf{x}_p$. (\emph{No requiere
        normalidad, aunque si puede asumirse la interpretación es más sencilla
        y pueden llevarse a cabo pruebas de significación.})
  \item Sean \textbf{combinaciones lineales} de
        $\mathbf{x}_1,\mathbf{x}_2,\dots,\mathbf{x}_p$, por lo que hay el
        mismo número de componentes que de variables.
  \item $\operatorname{var}(\mathbf{y}_1)\ge
        \operatorname{var}(\mathbf{y}_2)\ge\cdots\ge\operatorname{var}(\mathbf{y}_p)$.
  \item $\operatorname{corr}(\mathbf{y}_i,\mathbf{y}_j)=0\ \forall i\neq j$,
        por lo que no existen restricciones de linealidad entre componentes y
        la matriz de covarianzas de los componentes es diagonal. En la práctica,
        $\mat{S}$ y $\mat{R}$ son semidefinidas positivas, y serán definidas
        positivas cuando no exista colinealidad exacta.
  \item $\displaystyle\sum_{i=1}^p\operatorname{var}(\mathbf{y}_i)=
        \sum_{i=1}^p\operatorname{var}(\mathbf{x}_i)$.
\end{enumerate}

En resumen, el problema formal combina tres restricciones: ortogonalidad,
ordenación por varianza explicada e invariancia de la varianza total.

Suponemos variables centradas (media 0), de modo que
$\mathbf{X}$ tiene covarianza $\frac{1}{n}\mathbf{X}'\mathbf{X}$.
Se trata de encontrar un \textbf{subespacio de dimensión menor que $p$} tal
que, al proyectar sobre él los puntos, conserven su estructura con la menor
distorsión posible (distancias entre puntos originales y proyecciones lo más
pequeñas posible).

\subsection{Proyección y criterio de maximización}

Si consideramos un punto $\mathbf{x}_i$ y una dirección
$\mathbf{a}_1=(a_{11},\dots,a_{1p})'$ de norma unidad, la proyección de
$x_i$ sobre esta dirección es:
\[
  z_i = a_{11}x_{i1}+\cdots+a_{1p}x_{ip} = \mathbf{a}_1'\mathbf{x}_i,
\]
y el vector que representa esta proyección es $z_i\mathbf{a}_1$.
\begin{variablesprimera}
  \varitem{\mathbf{a}_1}{dirección del primer componente (vector de pesos)}
  \varitem{z_i}{proyección escalar de la observación $\mathbf{x}_i$ sobre $\mathbf{a}_1$}
  \varitem{r_i}{distancia entre la observación y su proyección en la dirección elegida}
\end{variablesprimera}

Por el teorema de Pitágoras: $\mathbf{x}_i'\mathbf{x}_i = z_i^2+r_i^2$,
y sumando para
todos los puntos:
\[
  \sum_{i=1}^n\mathbf{x}_i'\mathbf{x}_i = \sum_{i=1}^n z_i^2 + \sum_{i=1}^n r_i^2
  \quad\text{(constante)}.
\]
Por tanto:
\begin{mdframed}[style=resultado]
\[
  \min\sum_{i=1}^n r_i^2 \;\Longleftrightarrow\; \max\sum_{i=1}^n z_i^2.
\]
Como las $z_i$ tienen media 0 ($\bar{z}=\mathbf{a}_1'\bar{\mathbf{x}}=0$),
\textbf{maximizar la suma de cuadrados de las proyecciones equivale a
maximizar su varianza}.
\end{mdframed}

Este paso conecta explícitamente formulación geométrica (proyecciones) con
formulación estadística (varianza explicada).

% ══════════════════════════════════════════════════════════════════════════
\section{Cálculo de los componentes}
% ══════════════════════════════════════════════════════════════════════════

\subsection{Cálculo del primer componente}

Buscamos una combinación lineal $\mathbf{z}_1=\mathbf{X}\mathbf{a}_1$ con
\textbf{varianza máxima}. Como $\mathbf{z}_1$ tiene media 0:
\[
  \operatorname{var}(\mathbf{z}_1)
  = \frac{1}{n}\mathbf{z}_1'\mathbf{z}_1
  = \frac{1}{n}\mathbf{a}_1'\mathbf{X}'\mathbf{X}\mathbf{a}_1
  = \mathbf{a}_1'\mat{S}\mathbf{a}_1.
\]
Imponemos la restricción $\mathbf{a}_1'\mathbf{a}_1=1$ mediante un
multiplicador de Lagrange:
\[
  \mathcal{L} = \mathbf{a}_1'\mat{S}\mathbf{a}_1 - \lambda(\mathbf{a}_1'\mathbf{a}_1-1).
\]
\begin{variablesprimera}
  \varitem{\mat{S}}{matriz de covarianzas de las variables originales}
  \varitem{\mathcal{L}}{función de Lagrange del problema de optimización}
  \varitem{\lambda}{multiplicador de Lagrange que resulta ser un autovalor de $\mat{S}$}
\end{variablesprimera}
Derivando e igualando a cero:
\[
  \frac{\partial\mathcal{L}}{\partial\mathbf{a}_1} = 2\mat{S}\mathbf{a}_1-2\lambda\mathbf{a}_1=0
  \quad\Longrightarrow\quad \mat{S}\mathbf{a}_1=\lambda\mathbf{a}_1.
\]

\begin{mdframed}[style=resultado]
$\mathbf{a}_1$ es un \textbf{autovector} de la matriz $\mat{S}$ y $\lambda$
su correspondiente autovalor. Como
$\mathbf{a}_1'\mat{S}\mathbf{a}_1=\lambda\mathbf{a}_1'\mathbf{a}_1=\lambda$,
la varianza del componente es el autovalor: $\operatorname{var}(\mathbf{z}_1)=\lambda$.

Para maximizarla, $\lambda$ debe ser el \textbf{mayor valor propio de $\mat{S}$},
y $\mathbf{a}_1$ su vector asociado define los coeficientes de cada variable
en el primer componente principal.
\end{mdframed}

Este paso conecta de forma directa optimización y espectro de $\mat{S}$:
maximizar varianza bajo restricción de norma equivale a seleccionar el
autovector principal.

\begin{proposicion}[Caracterización variacional del primer componente]
El primer componente principal maximiza el cociente de Rayleigh:
\[
  \argmax_{\mathbf{a}\neq \mathbf{0}}
  \frac{\mathbf{a}'\mat{S}\mathbf{a}}{\mathbf{a}'\mathbf{a}}
  = \mathbf{a}_1,
\]
y el máximo alcanzado es el mayor autovalor de $\mat{S}$.
\end{proposicion}

\begin{proof}
Como $\mat{S}$ es simétrica, admite descomposición espectral
$\mat{S}=\mat{U}\boldsymbol{\Lambda}\mat{U}'$, con
$\lambda_1\ge\cdots\ge\lambda_p\ge 0$. Si
$\mathbf{a}=\mat{U}\mathbf{c}$ y $\|\mathbf{a}\|^2=\|\mathbf{c}\|^2=1$:
\[
  \mathbf{a}'\mat{S}\mathbf{a}
  = \mathbf{c}'\boldsymbol{\Lambda}\mathbf{c}
  = \sum_{j=1}^p \lambda_j c_j^2
  \le \lambda_1\sum_{j=1}^p c_j^2
  = \lambda_1.
\]
La cota se alcanza tomando $\mathbf{c}=(1,0,\dots,0)'$, es decir,
$\mathbf{a}=\mathbf{a}_1$ (autovector asociado a $\lambda_1$).
\end{proof}

\subsection{Cálculo del segundo componente}

El segundo componente:
\begin{enumerate}[label=\alph*)]
  \item Deberá ser \textbf{ortogonal} al primero: $\mathbf{a}_1'\mathbf{a}_2=0$.
  \item De entre las infinitas direcciones ortogonales posibles, escogemos la
        que extraiga el máximo de la varianza que queda sin explicar.
\end{enumerate}

La suma de las varianzas de $\mathbf{z}_1=\mathbf{X}\mathbf{a}_1$ y
$\mathbf{z}_2=\mathbf{X}\mathbf{a}_2$ debe ser máxima:
\[
  \mathcal{L} = \mathbf{a}_1'\mat{S}\mathbf{a}_1+\mathbf{a}_2'\mat{S}\mathbf{a}_2
  -\lambda_1(\mathbf{a}_1'\mathbf{a}_1-1)-\lambda_2(\mathbf{a}_2'\mathbf{a}_2-1).
\]
Derivando:
\[
  \mat{S}\mathbf{a}_1=\lambda_1\mathbf{a}_1, \qquad \mat{S}\mathbf{a}_2=\lambda_2\mathbf{a}_2.
\]
$\mathbf{a}_1$ y $\mathbf{a}_2$ son autovectores de $\mat{S}$. Tomando
vectores propios de norma uno: $\mathcal{L}=\lambda_1+\lambda_2$.

Como $\mathbf{a}_1'\mathbf{a}_2=0$:
$\operatorname{cov}(\mathbf{z}_1,\mathbf{z}_2)=\mathbf{a}_1'\mat{S}\mathbf{a}_2=0$,
con lo que $\mathbf{z}_1$ y $\mathbf{z}_2$ están \textbf{incorreladas}.
La construcción se repite inductivamente para componentes sucesivos: cada nueva
dirección maximiza varianza residual bajo ortogonalidad con las anteriores.

\subsection{Generalización}

\begin{proposicion}
Sea $h=\rg(\mat{S})\le \min(n,p)$. Existen $p$ componentes principales
algebraicamente definidos, pero solo $h$ tienen varianza estrictamente positiva.
Se obtienen calculando los autovalores
$\lambda_1\ge\lambda_2\ge\cdots\ge\lambda_p\ge0$ de la ecuación característica:
\[
  |\mat{S}-\lambda\mat{I}|=0, \qquad (\mat{S}-\lambda_i\mat{I})\mathbf{a}_i=0.
\]
Los $\lambda_i$ son reales y no negativos, al ser $\mat{S}$ simétrica
semidefinida positiva.
\end{proposicion}

Llamando $\mat{Z}$ a la matriz cuyas columnas son los valores de las $p$
componentes en los $n$ individuos, las nuevas variables se relacionan con las
originales mediante:
\begin{mdframed}[style=resultado]
\[
  \mat{Z} = \mathbf{X}\mat{A}, \qquad \mat{A}'\mat{A}=\mat{I},
\]
\end{mdframed}
donde $\mat{A}$ contiene en columnas los autovectores normalizados. Calcular
los componentes principales equivale a aplicar una \textbf{transformación
ortogonal} a las variables $\mathbf{X}$ (ejes originales) para obtener nuevas
variables $\mat{Z}$ incorreladas entre sí: se eligen nuevos ejes coordenados
que coincidan con los \emph{ejes naturales} de los datos.

\begin{proposicion}[Covarianza de las puntuaciones principales]
Si $\mat{Z}=\mat{X}_c\mat{A}$, con $\mat{A}'\mat{A}=\mat{I}$ y columnas de
$\mat{A}$ formadas por autovectores de $\mat{S}$, entonces:
\[
  \Cov(\mat{Z})=\mat{A}'\mat{S}\mat{A}=\boldsymbol{\Lambda},
\]
es decir, las componentes principales son incorreladas y sus varianzas son los
autovalores de $\mat{S}$.
\end{proposicion}

\begin{proof}
Como $\mat{X}_c$ está centrada, $\Cov(\mat{Z})=\frac{1}{n}\mat{Z}'\mat{Z}$.
Sustituyendo $\mat{Z}=\mat{X}_c\mat{A}$:
\[
  \Cov(\mat{Z})
  = \frac{1}{n}\mat{A}'\mat{X}_c'\mat{X}_c\mat{A}
  = \mat{A}'\mat{S}\mat{A}.
\]
Si $\mat{S}\mat{A}=\mat{A}\boldsymbol{\Lambda}$, entonces
$\mat{A}'\mat{S}\mat{A}=\mat{A}'\mat{A}\boldsymbol{\Lambda}
=\boldsymbol{\Lambda}$.
\end{proof}

% ══════════════════════════════════════════════════════════════════════════
\section{Propiedades de los componentes}
% ══════════════════════════════════════════════════════════════════════════

\begin{enumerate}
  \item \textbf{Conservación de la variabilidad total.}
        La suma de varianzas de los componentes es igual a la suma de varianzas
        de las variables originales, y la varianza generalizada de los
        componentes es igual a la original:
        \[
          \tr(\mat{S}) = \sum_{j=1}^p\Var(X_j)
          = \lambda_1+\cdots+\lambda_p
          = \sum_{j=1}^p\Var(Z_j).
        \]
        Si llamamos $\mat{S}_z$ a la matriz de covarianzas de los componentes
        (diagonal con términos $\lambda_i$):
        \[
          |\mat{S}_z| = \lambda_1\cdots\lambda_p
          = \prod_{j=1}^p\Var(Z_j).
        \]

  \item \textbf{Proporción de variabilidad explicada.}
        La proporción de variabilidad total explicada por el componente $h$ es:
        \[
          c_h = \frac{\lambda_h}{\sum_{i}\lambda_i}.
        \]

  \item \textbf{Covarianzas entre componentes y variables originales.}
        \[
          \operatorname{cov}(z_i;x_1,\dots,x_p) = \lambda_i\mathbf{a}_i
          = (\lambda_i a_{i1},\dots,\lambda_i a_{ip}).
        \]
        La matriz $p\times p$ de covarianzas entre componentes y variables
        originales es:
        \[
          \Cov(\mathbf{Z},\mathbf{X}) = \frac{1}{n}\mat{Z}'\mathbf{X}
          = \mat{A}'\mat{S} = \boldsymbol{\Lambda}\mat{A}',
        \]
        donde $\mat{A}$ contiene en columnas los vectores propios de $\mat{S}$
        y $\boldsymbol{\Lambda}$ es la matriz diagonal de valores propios.

  \item \textbf{Correlación entre un componente y una variable.}
        La correlación entre el componente $z_i$ y la variable $x_j$ es
        proporcional al coeficiente de esa variable en la definición del
        componente:
        \[
          r(z_i,x_j)
          = \frac{\operatorname{cov}(z_i,x_j)}{\sqrt{\operatorname{var}(z_i)\operatorname{var}(x_j)}}
          = \frac{\lambda_i a_{ij}}{\sqrt{\lambda_i s_j^2}}
          = a_{ij}\frac{\sqrt{\lambda_i}}{s_j}.
        \]
        El coeficiente de proporcionalidad es el cociente entre la desviación
        típica del componente y la de la variable.

  \item \textbf{Predicción óptima.}
        Las $r$ componentes principales ($r<p$) proporcionan la
        \textbf{predicción lineal óptima} del conjunto de variables $\mathbf{X}$.
\end{enumerate}

\begin{observacion}
La lectura conjunta de estas propiedades es clave: ACP no solo reduce
dimensión, sino que reorganiza la variabilidad en ejes ortogonales ordenados,
lo que facilita análisis, visualización e interpretación estadística.
\end{observacion}

% ══════════════════════════════════════════════════════════════════════════
\section{Análisis con correlaciones (análisis normado)}
% ══════════════════════════════════════════════════════════════════════════

Los componentes principales se obtienen maximizando la varianza de la
proyección. En términos de las variables originales:
\[
  \max M = \sum_{i=1}^p a_i^2 s_i^2
  + 2\sum_{i=1}^p\sum_{j=i+1}^p a_i a_j s_{ij}, \qquad \mathbf{a}'\mathbf{a}=1.
\]
Cuando las escalas de medida son muy distintas, la maximización dependerá de
estas escalas y las variables con valores más grandes tendrán más peso. La
solución es \textbf{estandarizar las variables antes de calcular los
componentes}, con lo que la expresión se reduce a:
\[
  1 + 2\sum_{i=1}^p\sum_{j=i+1}^p a_i a_j r_{ij},
\]
siendo $r_{ij}$ el coeficiente de correlación $\Rightarrow$ la solución
depende de las \emph{correlaciones}, no de las varianzas.

En términos prácticos:
\begin{itemize}
  \item ACP sobre $\mat{S}$ responde a la pregunta ``qué direcciones maximizan
        variabilidad absoluta''.
  \item ACP sobre $\mat{R}$ responde a ``qué direcciones maximizan estructura
        correlacional, eliminando efecto de escala''.
\end{itemize}

Así, la elección entre $\mat{S}$ y $\mat{R}$ no es un detalle técnico, sino
una decisión de modelización sobre qué noción de variabilidad se desea
priorizar.

\begin{mdframed}[style=resultado]
\begin{itemize}
  \item Si las variables originales están en \textbf{distintas unidades},
        conviene aplicar el análisis de la matriz de correlaciones
        (\emph{análisis normado}).
  \item Si tienen las mismas unidades, ambas alternativas son posibles.
  \item Si las diferencias de varianzas son informativas, no se deben
        estandarizar las variables.
\end{itemize}
\end{mdframed}

Los \textbf{componentes principales normados} se obtienen calculando los
vectores y valores propios de la matriz $\mat{R}$. Llamando $\lambda_i^R$ a
los autovalores de esa matriz (suponemos no singular):
\[
  \sum_{i=1}^p\lambda_i^R = \tr(\mat{R}) = p.
\]

\paragraph{Propiedades de los componentes extraídos de $\mat{R}$.}
\begin{itemize}
  \item La proporción de varianza explicada por la componente $j$ es:
        $\dfrac{\lambda_j^R}{p}$.
  \item Las correlaciones entre cada componente $z_j$ y las variables $X$
        originales vienen dadas por: $\mathbf{a}_j'\sqrt{\lambda_j}$,
        siendo $\mathbf{z}_j=\mathbf{X}\mathbf{a}_j$.
\end{itemize}

% ══════════════════════════════════════════════════════════════════════════
\section{Matriz de saturaciones o de cargas factoriales}
% ══════════════════════════════════════════════════════════════════════════

\begin{definicion}[Matriz de saturaciones]
La \textbf{matriz de saturaciones} $\mat{F}$ muestra las correlaciones entre
las variables originales y los componentes principales. Se calcula como:
\begin{mdframed}[style=resultado]
\[
  \mat{F} = \mat{A}\boldsymbol{\Lambda}^{1/2},
\]
\end{mdframed}
donde $\mat{A}$ es la matriz de autovectores normalizados de la matriz
analizada y $\boldsymbol{\Lambda}$ la matriz diagonal de autovalores.
\end{definicion}

Interpretativamente, cada carga factorial es una correlación variable--
componente: magnitud alta implica fuerte asociación de la variable con el eje
latente correspondiente.

\begin{itemize}
  \item Cada \textbf{fila} representa una variable original.
  \item Cada \textbf{columna} representa un componente principal.
  \item Cuanto mayor es el valor absoluto de una saturación para un
        componente, más fuerte es la relación entre esa variable y el
        componente.
  \item Se interpreta como la \textbf{matriz del cambio de la base original}
        a la base constituida por los componentes.
  \item Expresada en puntuaciones típicas, las saturaciones son las
        correlaciones entre cada $X_i$ e $Y_i$:
        $\mat{D}^{-1/2}\mat{A}\boldsymbol{\Lambda}^{1/2} = \mat{R}_{xy}$,
        siendo $\mat{D}=\diag(s_1^2,\dots,s_p^2)$ la matriz diagonal de
        varianzas de las variables originales.
\end{itemize}

\begin{observacion}
Las cargas factoriales son el puente entre formulación algebraica e
interpretación sustantiva: permiten nombrar componentes a partir del patrón de
variables que más contribuye a cada eje.
\end{observacion}

% ══════════════════════════════════════════════════════════════════════════
\section{Resumen: cálculo de los componentes principales}
% ══════════════════════════════════════════════════════════════════════════

\subsection*{A partir de la matriz $\mat{S}$}

\begin{enumerate}[label=\Roman*)]
  \item Cálculo de la matriz de covarianzas:
        $\mat{S} = \dfrac{\mat{X}_c'\mat{X}_c}{n-1}$, donde
        $\mat{X}_c$ es la matriz de datos centrados.
  \item Cálculo de la matriz $\boldsymbol{\Lambda}$ de autovalores de $\mat{S}$:
        $|\mat{S}-\lambda\mat{I}|=0$.
  \item Cálculo de la matriz $\mat{A}$ de autovectores:
        $(\mat{S}-\lambda_i\mat{I})\mathbf{a}_i=0$
        (autovectores normalizados de $\mat{S}$; el primero corresponde a la
        mayor raíz latente).
  \item Cálculo de la matriz $\mat{F}$ de saturaciones:
        $\mat{F}=\mat{A}\boldsymbol{\Lambda}^{1/2}$.
  \item Cálculo de la matriz de puntuaciones de los sujetos en los
        componentes: $\mat{Y}=\mat{X}_c\mat{A}$.
  \item Cálculo de la varianza de cada componente:
        $\operatorname{var}(Y_i)=\mathbf{a}_i'\mat{S}\mathbf{a}_i=\lambda_i$.
  \item Cálculo de la proporción de varianza explicada:
        $c_i = \lambda_i/\sum\lambda_i$.
\end{enumerate}

\subsection*{A partir de la matriz $\mat{R}$}

\begin{enumerate}[label=\Roman*)]
  \item Cálculo de la matriz de correlaciones:
        $\mat{R} = \dfrac{\mat{Z}_x'\mat{Z}_x}{n-1}$, con $\mat{Z}_x$ datos
        tipificados.
  \item Cálculo de la matriz $\boldsymbol{\Lambda}$ de autovalores de $\mat{R}$:
        $|\mat{R}-\lambda\mat{I}|=0$.
  \item Cálculo de la matriz $\mat{A}$ de autovectores:
        $(\mat{R}-\lambda_i\mat{I})\mathbf{a}=0$.
  \item Cálculo de la matriz $\mat{F}$ de saturaciones:
        $\mat{F}=\mat{R}_{\mathbf{z}_x,\mathbf{z}_y}=\mat{A}\boldsymbol{\Lambda}^{1/2}$.
  \item Cálculo de la matriz de puntuaciones típicas:
        $\mat{Z}_y=\mat{Z}_x\mat{A}$.
  \item Cálculo de la varianza de cada componente:
        $\operatorname{var}(\mathbf{z}_{y_i})=\mat{A}'\mat{R}\mat{A}=\boldsymbol{\Lambda}
        =\mathbf{a}_i'\mat{R}\mathbf{a}_i=\lambda_i$;
        $\operatorname{cov}(\mathbf{z}_x,\mathbf{z}_y)=\mat{A}\boldsymbol{\Lambda}$.
  \item Cálculo de la proporción de varianza explicada:
        $c_i=\lambda_i/p$.
\end{enumerate}

% ══════════════════════════════════════════════════════════════════════════
\section{Interpretación de los componentes}
% ══════════════════════════════════════════════════════════════════════════

\subsection{Interpretación geométrica}

Las combinaciones lineales $\mat{Y}=\mathbf{X}\mat{A}$ representan un nuevo
sistema de coordenadas obtenido mediante una \textbf{rotación ortogonal a
través del origen}:
\[
  X_1, X_2, \dots, X_p
  \;\xrightarrow{\text{rotación ortogonal}}\;
  Y_1, Y_2, \dots, Y_p.
\]
La rotación es ortogonal porque $\mat{A}$ es una matriz ortogonal que, para
dos componentes, puede expresarse como:
\[
  \mat{A} = \begin{pmatrix}\cos\theta & -\sin\theta\\
             \sin\theta & \cos\theta\end{pmatrix}.
\]
La matriz $\mat{A}$ rota los ejes $X_1, X_2, \dots, X_p$ un ángulo $\theta$
de manera que el primer componente reúna el máximo de la varianza contenida
en los datos (y el segundo la restante).

\begin{ejemplo}
\[
  \mat{A}=\begin{pmatrix}0{,}722&-0{,}692\\0{,}692&0{,}722\end{pmatrix},
  \quad\cos\theta=0{,}722,\quad\sin\theta=0{,}692
  \quad\Rightarrow\quad\theta=43{,}78^\circ.
\]
Los ejes se han girado $43{,}78^\circ$ de manera que el componente $y_1$ tiene
gran variabilidad y el componente $y_2$ muy pequeña.
\end{ejemplo}

\subsection{Interpretación sustantiva}

\begin{itemize}
  \item \textbf{Alta correlación positiva} entre todas las variables
        $\Rightarrow$ primer componente principal con todas sus coordenadas
        del mismo signo. Puede interpretarse como un \emph{promedio ponderado}
        de todas las variables $\Rightarrow$ factor global de ``tamaño''.
  \item Las \textbf{restantes componentes} se interpretan como factores ``de
        forma'', con coordenadas positivas y negativas que contraponen unos
        grupos de variables frente a otros.
\end{itemize}

% ══════════════════════════════════════════════════════════════════════════
\section{Número de factores a extraer}
% ══════════════════════════════════════════════════════════════════════════

No se ha desarrollado una base cuantitativa exacta para decidir el número de
componentes a extraer. Los criterios más habituales son:

\begin{enumerate}[label=\Roman*)]
  \item \textbf{Criterio de raíz latente (Kaiser-Guttman, 1954):} cualquier
        factor individual debería justificar la varianza de al menos una
        variable. Se consideran significativos todos los factores con
        autovalores $\ge1$ (los que explican menos de una variable se
        desestiman).
        \begin{itemize}
          \item Es más fiable cuando $20\le p\le50$.
          \item $p<20$: criterio conservador (demasiados pocos factores).
          \item $p>50$: demasiados factores.
        \end{itemize}
        \emph{Es una regla arbitraria: una variable independiente del resto
        puede tener autovalor $>1$, pero si está incorrelada con el resto puede
        ser poco relevante para el análisis.}

  \item \textbf{Criterio a priori:} se fija el número de factores antes de
        realizar el análisis.

  \item \textbf{Criterio del porcentaje de la varianza:} se basa en obtener
        un porcentaje acumulado especificado de la varianza total.
        \begin{itemize}
          \item Ciencias naturales: $\ge95\%$ de la varianza total.
          \item Ciencias sociales: $\ge60\%$ de la varianza total.
        \end{itemize}
        \emph{La regla es de nuevo arbitraria y debe aplicarse con cuidado:
        pueden existir componentes fuera de esos porcentajes muy adecuados para
        explicar la ``forma'' de las variables.}

  \item \textbf{Gráfico de desmoronamiento o \emph{scree plot}:} gráfico de
        $\lambda_i$ frente a $i$. Las componentes se seleccionan hasta que las
        restantes tengan aproximadamente el mismo valor de $\lambda_i$
        (``codo'' en el gráfico). A partir de ese punto los autovalores no
        aportan información adicional relevante.
\end{enumerate}

\begin{observacion}
En práctica, la selección final combina criterios numéricos y criterio
interpretativo: una solución útil debe ser simultáneamente parsimoniosa y
sustantivamente legible.
\end{observacion}

% ══════════════════════════════════════════════════════════════════════════
\section{Rotaciones}
% ══════════════════════════════════════════════════════════════════════════

Extraer componentes para que tengan varianza máxima puede conducir a
componentes muy complejas (correlación con la mayor parte de las variables
originales), lo que dificulta la interpretación. Para facilitarla, se pueden
\textbf{transformar los componentes mediante una rotación} (cambio de base):
\begin{itemize}
  \item Se utiliza sobre todo en \textbf{análisis factorial}.
  \item La determinación del ángulo de rotación de las $p$ componentes
        obtenidas se lleva a cabo con criterios matemáticos de optimización.
\end{itemize}

\begin{definicion}[Rotación Varimax (Kaiser, 1958)]
El método más utilizado. Es una \textbf{rotación ortogonal} que simplifica
por columnas maximizando la varianza de las saturaciones de cada componente.
\end{definicion}

El resultado es que las columnas pasan a tener \textbf{patrones de
saturaciones muy dispersos}:
\begin{itemize}
  \item Cargas elevadas en un solo factor y bajas en los otros.
  \item Mejor identificación de las variables más relacionadas con cada
        factor, mejorando la interpretación y la claridad de los resultados.
\end{itemize}

\begin{observacion}
Otras rotaciones ortogonales: Quartimax (simplifica filas), Equimax
(compromiso). Rotaciones oblicuas (factores correlados): Oblimin, Promax.
Las rotaciones no alteran la varianza total explicada, solo redistribuyen
la varianza entre los componentes.
\end{observacion}

% ══════════════════════════════════════════════════════════════════════════
%  TEMA 4
% ══════════════════════════════════════════════════════════════════════════
\iniciotema{4}{Análisis Discriminante}{Análisis Discriminante}

% ══════════════════════════════════════════════════════════════════════════
\section{Introducción}
% ══════════════════════════════════════════════════════════════════════════

El \textbf{Análisis Discriminante Lineal} (ADL o \emph{LDA}) es una técnica
de \textbf{clasificación supervisada}: los grupos están definidos a priori y
se construye una regla para asignar nuevas observaciones al grupo más probable
según sus variables explicativas.

Desde un punto de vista metodológico, el problema central es construir una
regla de decisión con buen equilibrio entre separabilidad geométrica de clases
y riesgo de clasificación.

\begin{itemize}
  \item Es un método clásico de \textbf{reconocimiento de patrones}.
  \item Los grupos deben ser \textbf{mutuamente excluyentes} y, en el esquema
        estándar, \textbf{exhaustivos}.
  \item Es apropiado cuando la variable respuesta es \textbf{cualitativa}
        (grupos) y las variables predictoras son \textbf{cuantitativas}.
  \item Se apoya en una combinación lineal de predictores, análoga al valor
        teórico ya estudiado en temas previos.
  \item Se necesita una \textbf{muestra de entrenamiento} (grupo conocido) y,
        para evaluar el rendimiento, una \textbf{muestra de validación}
        (\emph{test sample}) o validación cruzada.
  \item En variables de respuesta ordinales puede usarse un \emph{enfoque de
        extremos polares} (comparar solo grupos extremos), útil en algunos
        contextos aunque con pérdida de información intermedia.
\end{itemize}

\begin{mdframed}[style=resultado]
\textbf{Hipótesis canónicas del LDA clásico:}
\begin{enumerate}
  \item Normalidad multivariante dentro de cada clase.
  \item Matriz de covarianzas común entre clases.
  \item Independencia entre observaciones.
  \item Colinealidad no exacta entre predictores.
\end{enumerate}
\end{mdframed}

\begin{definicion}[Función discriminante lineal]
Para una observación $k$, una función discriminante lineal puede escribirse
como:
\begin{mdframed}[style=resultado]
\[
  z_k = a + w_1x_{1k} + w_2x_{2k} + \cdots + w_px_{pk}.
\]
\end{mdframed}
\begin{variablesprimera}
  \varitem{z_k}{puntuación discriminante de la observación $k$}
  \varitem{a}{término independiente de la función discriminante}
  \varitem{x_{jk}}{valor de la variable $j$ en la observación $k$}
\end{variablesprimera}
Los pesos $w_i$ se eligen para \textbf{maximizar la separación entre grupos}
frente a la variabilidad dentro de los grupos.
\end{definicion}

Cada grupo tiene un \textbf{centroide} (media de puntuaciones discriminantes).
Cuanto más separados estén los centroides, mayor capacidad discriminante.

\begin{nota}
El enfoque clásico se debe a \textbf{Fisher (1936)} y es óptimo bajo
normalidad multivariante con matriz de covarianzas común en los grupos.
Cuando estos supuestos fallan, se suelen considerar alternativas como
QDA, regresión logística multinomial, árboles, k-NN o SVM.
\end{nota}

% ══════════════════════════════════════════════════════════════════════════
\section{Caso de dos poblaciones}
% ══════════════════════════════════════════════════════════════════════════

\subsection{Enfoque bayesiano general}

Sean dos poblaciones $\mathcal{P}_1,\mathcal{P}_2$ con densidades
$f_1(\mathbf{x})$, $f_2(\mathbf{x})$ y probabilidades a priori
$\pi_1,\pi_2$ ($\pi_1+\pi_2=1$). Para un nuevo elemento $\mathbf{x}_0$:

\begin{mdframed}[style=resultado]
\[
  P(\mathcal{P}_i\mid \mathbf{x}_0)=
  \frac{\pi_i f_i(\mathbf{x}_0)}{\pi_1f_1(\mathbf{x}_0)+\pi_2f_2(\mathbf{x}_0)},
  \qquad i=1,2.
\]
\end{mdframed}
\begin{variablesprimera}
  \varitem{\mathcal{P}_i}{población o clase $i$}
  \varitem{\pi_i}{probabilidad a priori de la clase $i$}
  \varitem{f_i(\mathbf{x})}{densidad de $\mathbf{x}$ en la clase $i$}
  \varitem{\mathbf{x}_0}{nueva observación a clasificar}
  \varitem{P(\mathcal{P}_i\mid\mathbf{x}_0)}{probabilidad a posteriori de pertenecer a la clase $i$}
\end{variablesprimera}

Regla bayesiana básica: asignar $\mathbf{x}_0$ al grupo con mayor probabilidad
a posteriori; equivalentemente, al grupo que maximiza
$\pi_i f_i(\mathbf{x}_0)$. Esta equivalencia permite pasar de una regla
probabilística a una regla operativa de decisión.

La secuencia conceptual es: hipótesis probabilística sobre clases, cálculo de
posteriores y, finalmente, regla de asignación por minimización de riesgo.

\subsubsection{Regla con costes de clasificación}

Si el coste de clasificar en 2 algo de 1 es $c(2|1)$ y el coste inverso es
$c(1|2)$:
\begin{definicion}[Regla de decisión]
Una regla de decisión particiona el espacio muestral
$\mathcal{E}_{\mathbf{x}}$ en dos regiones $\mathcal{A}_1,\mathcal{A}_2$:
\[
  \mathcal{A}_2=\mathcal{E}_{\mathbf{x}}\setminus\mathcal{A}_1,\qquad
  \mathbf{x}_0\in\mathcal{A}_i \Rightarrow \text{asignar a } \mathcal{P}_i.
\]
Si se conocen costes de clasificación, se asigna a la clase con
\textbf{menor riesgo condicional esperado}.
\end{definicion}

\[
  R(d_1\mid \mathbf{x}_0)=c(1|2)\,P(\mathcal{P}_2\mid\mathbf{x}_0),\qquad
  R(d_2\mid \mathbf{x}_0)=c(2|1)\,P(\mathcal{P}_1\mid\mathbf{x}_0).
\]
Se asigna a $\mathcal{P}_2$ si $R(d_2\mid\mathbf{x}_0)<R(d_1\mid\mathbf{x}_0)$,
es decir:
\begin{mdframed}[style=resultado]
\[
  c(1|2)\,\pi_2f_2(\mathbf{x}_0) > c(2|1)\,\pi_1f_1(\mathbf{x}_0).
\]
\end{mdframed}

\begin{itemize}
  \item Si $c(1|2)=c(2|1)$, la decisión depende de $\pi_i f_i(\mathbf{x}_0)$.
  \item Si además $\pi_1=\pi_2$, la regla se reduce a comparar densidades:
        clasificar en la población con mayor verosimilitud en $\mathbf{x}_0$.
\end{itemize}

Por tanto, priori y costes no son correcciones accesorias: desplazan
explícitamente la frontera de decisión.

\subsection{Enfoque bayesiano en poblaciones normales}

Si $\mathbf{X}\mid \mathcal{P}_i\sim\mathcal{N}_p(\boldsymbol{\mu}_i,\boldsymbol{\Sigma})$
con \textbf{misma} matriz $\boldsymbol{\Sigma}$:
\[
  f_i(\mathbf{x})=
  \frac{1}{(2\pi)^{p/2}|\boldsymbol{\Sigma}|^{1/2}}
  \exp\!\left\{-\frac12(\mathbf{x}-\boldsymbol{\mu}_i)'
  \boldsymbol{\Sigma}^{-1}(\mathbf{x}-\boldsymbol{\mu}_i)\right\}.
\]

Bajo normalidad multivariante y covarianza común, la frontera de decisión se
vuelve lineal en $\mathbf{x}$ usando las funciones:
\begin{mdframed}[style=resultado]
\[
  \delta_i(\mathbf{x})=
  \mathbf{x}'\boldsymbol{\Sigma}^{-1}\boldsymbol{\mu}_i
  -\frac12\boldsymbol{\mu}_i'\boldsymbol{\Sigma}^{-1}\boldsymbol{\mu}_i
  +\log\pi_i,
  \qquad i=1,2.
\]
\end{mdframed}
\begin{proof}
Partimos de la regla bayesiana
$\argmax_i\{\pi_i f_i(\mathbf{x})\}$. Tomando logaritmos y sustituyendo la
densidad normal multivariante con covarianza común:
\[
  \log\pi_i-\frac{1}{2}(\mathbf{x}-\boldsymbol{\mu}_i)'
  \boldsymbol{\Sigma}^{-1}(\mathbf{x}-\boldsymbol{\mu}_i)
  -\frac{1}{2}\log|\boldsymbol{\Sigma}|-\frac{p}{2}\log(2\pi).
\]
Los dos últimos términos no dependen de $i$, por lo que no afectan a la
decisión. Al expandir el término cuadrático:
\[
  -\frac12\mathbf{x}'\boldsymbol{\Sigma}^{-1}\mathbf{x}
  +\mathbf{x}'\boldsymbol{\Sigma}^{-1}\boldsymbol{\mu}_i
  -\frac12\boldsymbol{\mu}_i'\boldsymbol{\Sigma}^{-1}\boldsymbol{\mu}_i
  +\log\pi_i.
\]
El término $-\frac12\mathbf{x}'\boldsymbol{\Sigma}^{-1}\mathbf{x}$ también es
común a todas las clases. Eliminándolo se obtiene exactamente
$\delta_i(\mathbf{x})$.
\end{proof}
\begin{variablesprimera}
  \varitem{\delta_i(\mathbf{x})}{función discriminante lineal de la clase $i$}
  \varitem{\boldsymbol{\mu}_i}{vector de medias de la clase $i$}
  \varitem{\log\pi_i}{ajuste por probabilidad a priori}
\end{variablesprimera}
Se asigna al grupo con mayor $\delta_i(\mathbf{x})$.

Si $\pi_1=\pi_2$ y costes iguales, equivale a clasificar por \textbf{distancia
de Mahalanobis mínima} al centro:
\[
  D_i^2(\mathbf{x})=(\mathbf{x}-\boldsymbol{\mu}_i)'
  \boldsymbol{\Sigma}^{-1}
  (\mathbf{x}-\boldsymbol{\mu}_i),\qquad
  \text{elegir } \arg\min_i D_i^2(\mathbf{x}).
\]
\begin{variablesprimera}
  \varitem{D_i^2(\mathbf{x})}{distancia de Mahalanobis de $\mathbf{x}$ a la clase $i$}
\end{variablesprimera}
Si $\boldsymbol{\Sigma}=\sigma^2\mat{I}$, la regla se reduce a distancia
euclídea.

\begin{mdframed}[style=resultado]
Con normalidad multivariante y matriz común, la regla con costes y priori
puede escribirse como:
\[
  D_1^2(\mathbf{x})-D_2^2(\mathbf{x})
  > 2\log\!\left(\frac{c(2|1)\pi_1}{c(1|2)\pi_2}\right)
  \;\Longrightarrow\;
  \text{clasificar en } \mathcal{P}_2.
\]
Esta expresión evidencia que el umbral de decisión se desplaza según costes y
priori.
\end{mdframed}

\subsubsection{Probabilidad a posteriori en función de distancias}

Para dos grupos:
\begin{mdframed}[style=resultado]
\[
  P(\mathcal{P}_1\mid\mathbf{x})=
  \frac{1}{
    1+\dfrac{\pi_2}{\pi_1}\exp\!\left[-\frac12\left(D_2^2(\mathbf{x})-D_1^2(\mathbf{x})\right)\right]
  }.
\]
\end{mdframed}
Cuanto menor sea $D_1^2$ frente a $D_2^2$, mayor será la probabilidad
a posteriori de pertenecer a $\mathcal{P}_1$.

\subsubsection{Estimación de parámetros en la práctica}

En aplicaciones reales, $\pi_g$, $\boldsymbol{\mu}_g$ y
$\boldsymbol{\Sigma}$ no son conocidos y se estiman en la muestra de
entrenamiento:
\[
  \hat{\pi}_g = \frac{n_g}{n},\qquad
  \hat{\boldsymbol{\mu}}_g = \frac{1}{n_g}\sum_{i:y_i=g}\mathbf{x}_i,
\]
\[
  \hat{\boldsymbol{\Sigma}}
  = \frac{1}{n-G}
    \sum_{g=1}^{G}\sum_{i:y_i=g}
    (\mathbf{x}_i-\hat{\boldsymbol{\mu}}_g)
    (\mathbf{x}_i-\hat{\boldsymbol{\mu}}_g)'.
\]

La regla empírica queda:
\[
  \hat{\delta}_g(\mathbf{x})=
  \mathbf{x}'\hat{\boldsymbol{\Sigma}}^{-1}\hat{\boldsymbol{\mu}}_g
  -\frac12\hat{\boldsymbol{\mu}}_g'\hat{\boldsymbol{\Sigma}}^{-1}\hat{\boldsymbol{\mu}}_g
  +\log\hat{\pi}_g,
\]
asignando al grupo con mayor $\hat{\delta}_g(\mathbf{x})$.

\subsubsection{Ejemplo numérico (dos pintores)}

Se observan dos variables: profundidad de trazo y proporción del lienzo.
Con medias:
\[
  \boldsymbol{\mu}_A=(2,\,0.8)',\qquad
  \boldsymbol{\mu}_B=(2.3,\,0.7)',
\]
y matriz común
\[
  \boldsymbol{\Sigma}=
  \begin{pmatrix}
    0.25 & 0.025\\
    0.025 & 0.01
  \end{pmatrix},
  \quad
  \boldsymbol{\Sigma}^{-1}\approx
  \begin{pmatrix}
    5.33 & -13.33\\
    -13.33 & 133.33
  \end{pmatrix}.
\]

Para $\mathbf{x}_0=(2.1,\,0.75)'$:
\[
  D_A^2(\mathbf{x}_0)\approx0.520,\qquad
  D_B^2(\mathbf{x}_0)\approx0.813.
\]
Como $D_A^2<D_B^2$, se clasifica en el grupo $A$. El orden lógico del cálculo
es: (i) estimar/usar $\boldsymbol{\Sigma}^{-1}$, (ii) evaluar ambas distancias
cuadráticas, y (iii) aplicar la regla de mínima distancia bajo hipótesis de
covarianza común.
El ejemplo ilustra que la regla no se basa en distancia euclídea, sino en
distancia reescalada por la estructura de covarianza común.

Si $\pi_A=\pi_B$:
\[
  P(A\mid\mathbf{x}_0)\approx
  \frac{1}{1+\exp\!\left[-\frac12(0.813-0.520)\right]}
  \approx 0.537.
\]
La asignación es razonable, pero con incertidumbre apreciable.
Interpretativamente, el punto está más cerca de $A$ en geometría de Mahalanobis
pero no muy alejado de la frontera, de ahí la posterior moderada.

\subsubsection{Probabilidad teórica de error (caso simétrico)}

En dos grupos normales con covarianza común, y con costes/priori iguales, si
\[
  \Delta^2=
  (\boldsymbol{\mu}_1-\boldsymbol{\mu}_2)'
  \boldsymbol{\Sigma}^{-1}
  (\boldsymbol{\mu}_1-\boldsymbol{\mu}_2),
\]
la probabilidad de clasificación errónea en cada sentido es:
\begin{mdframed}[style=resultado]
\[
  P(2|1)=P(1|2)=1-\Phi\!\left(\frac{\Delta}{2}\right).
\]
\end{mdframed}

En el ejemplo, $\Delta\approx1.6166$, por lo que
$1-\Phi(\Delta/2)\approx 0.2095$ (error teórico alrededor del 21\%).

% ══════════════════════════════════════════════════════════════════════════
\section{Generalización para varias poblaciones normales}
% ══════════════════════════════════════════════════════════════════════════

Para $G$ poblaciones, la región de decisión de la clase $g$ es:
\[
  \mathcal{A}_g=
  \left\{
    \mathbf{x}\in\R^p:\;
    \pi_g f_g(\mathbf{x})>\pi_i f_i(\mathbf{x}),\ \forall i\neq g
  \right\}.
\]
\begin{variablesprimera}
  \varitem{G}{número total de clases}
  \varitem{g}{índice de clase}
  \varitem{\mathcal{A}_g}{región de decisión asignada a la clase $g$}
\end{variablesprimera}

\begin{itemize}
  \item Esta formulación compara simultáneamente todas las clases.
  \item También puede implementarse por comparaciones dos a dos entre
        poblaciones, pero la regla global anterior es la forma canónica.
\end{itemize}

Esta extensión preserva la misma lógica del caso binario: maximizar
probabilidad a posteriori (o minimizar riesgo) en todo el conjunto de clases.

Si $f_g$ son normales con \textbf{covarianza común} $\boldsymbol{\Sigma}$:
\[
  \mathbf{X}\mid g\sim\mathcal{N}_p(\boldsymbol{\mu}_g,\boldsymbol{\Sigma}),
\]
la regla es asignar al grupo con mayor:
\[
  \delta_g(\mathbf{x})=
  \mathbf{x}'\boldsymbol{\Sigma}^{-1}\boldsymbol{\mu}_g
  -\frac12\boldsymbol{\mu}_g'\boldsymbol{\Sigma}^{-1}\boldsymbol{\mu}_g
  +\log\pi_g.
\]
Con priori iguales, equivale a minimizar $D_g^2(\mathbf{x})$.

\subsection{Fronteras de decisión}

La frontera entre grupos $i,j$ se obtiene de
$\delta_i(\mathbf{x})=\delta_j(\mathbf{x})$, es decir, un \textbf{hiperplano}:
\[
  (\boldsymbol{\mu}_i-\boldsymbol{\mu}_j)'\boldsymbol{\Sigma}^{-1}\mathbf{x}
  -\frac12
  \left(
    \boldsymbol{\mu}_i'\boldsymbol{\Sigma}^{-1}\boldsymbol{\mu}_i
    -\boldsymbol{\mu}_j'\boldsymbol{\Sigma}^{-1}\boldsymbol{\mu}_j
  \right)
  +\log\!\frac{\pi_i}{\pi_j}=0.
\]

\begin{proof}
La igualdad $\delta_i(\mathbf{x})=\delta_j(\mathbf{x})$ elimina los términos
comunes de segundo grado en $\mathbf{x}$ (procedentes de
$-\tfrac12\mathbf{x}'\boldsymbol{\Sigma}^{-1}\mathbf{x}$), ya que la
covarianza es compartida por todas las clases. Por ello, el contorno de
indiferencia queda descrito por una ecuación afín en $\mathbf{x}$, es decir,
un hiperplano.
\end{proof}

% ══════════════════════════════════════════════════════════════════════════
\section{Función lineal de Fisher}
% ══════════════════════════════════════════════════════════════════════════

\subsection{Caso de dos grupos}

El criterio de Fisher puede leerse como una reformulación geométrica del
objetivo bayesiano gaussiano: hallar una proyección escalar
$z=\boldsymbol{\alpha}'\mathbf{x}$ que maximice separación intergrupos
relativa a dispersión intragrupo.

Así, el paso desde el enfoque bayesiano al de Fisher no es un cambio de
problema, sino un cambio de formulación: de regla probabilística a criterio
variacional geométrico.

Con medias muestrales $\bar{\mathbf{x}}_1,\bar{\mathbf{x}}_2$ y matriz
intra-grupos agrupada $\mat{S}_w$, el criterio es:
\begin{mdframed}[style=resultado]
\[
  J(\boldsymbol{\alpha})=
  \frac{\left[\boldsymbol{\alpha}'(\bar{\mathbf{x}}_2-\bar{\mathbf{x}}_1)\right]^2}
       {\boldsymbol{\alpha}'\mat{S}_w\boldsymbol{\alpha}},
  \qquad
  \boldsymbol{\alpha}\propto
  \mat{S}_w^{-1}(\bar{\mathbf{x}}_2-\bar{\mathbf{x}}_1).
\]
\end{mdframed}
\begin{variablesprimera}
  \varitem{\boldsymbol{\alpha}}{dirección de proyección discriminante de Fisher}
  \varitem{J(\boldsymbol{\alpha})}{cociente entre separación intergrupos y dispersión intragrupo}
  \varitem{\mat{S}_w}{matriz de dispersión intra-grupos agrupada}
\end{variablesprimera}

\begin{proof}
Definamos $\mathbf{d}=\bar{\mathbf{x}}_2-\bar{\mathbf{x}}_1$. Maximizar
$J(\boldsymbol{\alpha})$ equivale a maximizar
$\left(\boldsymbol{\alpha}'\mathbf{d}\right)^2$ bajo la restricción
$\boldsymbol{\alpha}'\mat{S}_w\boldsymbol{\alpha}=1$.
Con multiplicadores de Lagrange:
\[
  \mathcal{L}=
  \left(\boldsymbol{\alpha}'\mathbf{d}\right)^2
  -\lambda\!\left(\boldsymbol{\alpha}'\mat{S}_w\boldsymbol{\alpha}-1\right).
\]
La condición de primer orden conduce a
$\mathbf{d}\mathbf{d}'\boldsymbol{\alpha}=\lambda\mat{S}_w\boldsymbol{\alpha}$,
cuyo vector propio no nulo es proporcional a $\mat{S}_w^{-1}\mathbf{d}$.
Por tanto,
$\boldsymbol{\alpha}\propto \mat{S}_w^{-1}(\bar{\mathbf{x}}_2-\bar{\mathbf{x}}_1)$.
\end{proof}

Así, la dirección discriminante coincide con la dirección que maximiza la
distancia de Mahalanobis entre los centroides proyectados.

\begin{itemize}
  \item Si $\mat{S}_w=\mat{I}$, la dirección es paralela al vector que une
        medias (distancia euclídea).
  \item Con costes y priori iguales, en el eje proyectado se usa como umbral
        el punto medio entre medias proyectadas.
\end{itemize}

\subsection{Caso de varios grupos}

Para $G$ grupos se buscan
$r=\min(p,G-1)$ variables canónicas:
\[
  z_j=\mathbf{u}_j'\mathbf{x},\qquad j=1,\dots,r,
\]
que maximizan separación entre grupos y minimizan dispersión intragrupos.

Definimos:
\[
  \mat{W}=\sum_{g=1}^{G}\sum_{i=1}^{n_g}
  (\mathbf{x}_{ig}-\bar{\mathbf{x}}_g)(\mathbf{x}_{ig}-\bar{\mathbf{x}}_g)',
\]
\[
  \mat{B}=\sum_{g=1}^{G}
  n_g(\bar{\mathbf{x}}_g-\bar{\mathbf{x}})
  (\bar{\mathbf{x}}_g-\bar{\mathbf{x}})'.
\]
\begin{variablesprimera}
  \varitem{\mat{W}}{matriz de dispersión intra-grupos}
  \varitem{\mat{B}}{matriz de dispersión entre grupos}
  \varitem{n_g}{tamaño muestral del grupo $g$}
  \varitem{\mathbf{u}}{dirección canónica candidata}
\end{variablesprimera}

Cada dirección se obtiene maximizando:
\[
  \phi(\mathbf{u})=
  \frac{\mathbf{u}'\mat{B}\mathbf{u}}{\mathbf{u}'\mat{W}\mathbf{u}},
\]
lo que conduce al problema de autovalores:
\begin{mdframed}[style=resultado]
\[
  \mat{B}\mathbf{u}=\lambda\mat{W}\mathbf{u}
  \quad\Longleftrightarrow\quad
  \mat{W}^{-1}\mat{B}\mathbf{u}=\lambda\mathbf{u}.
\]
\end{mdframed}
\begin{variablesprimera}
  \varitem{\lambda}{autovalor asociado a una dirección canónica}
\end{variablesprimera}

Los autovectores asociados a los mayores autovalores proporcionan
las variables canónicas más discriminantes.

% ══════════════════════════════════════════════════════════════════════════
\section{Variables canónicas discriminantes}
% ══════════════════════════════════════════════════════════════════════════

Sea $\mat{U}_r=[\mathbf{u}_1,\dots,\mathbf{u}_r]$.
Para una observación $\mathbf{x}$:
\[
  \mathbf{z}=\mat{U}_r'\mathbf{x}\in\R^r.
\]

\begin{itemize}
  \item $\mathbf{z}$ son las coordenadas en el espacio discriminante.
  \item Cada grupo tiene centroide canónico
        $\bar{\mathbf{z}}_g=\mat{U}_r'\bar{\mathbf{x}}_g$.
  \item Una regla práctica de clasificación es asignar $\mathbf{z}_0$
        al centroide más próximo (normalmente distancia euclídea en ese espacio).
\end{itemize}

Geométricamente, la transformación canónica busca un espacio de dimensión baja
donde los centroides queden lo más separados posible con métrica intragrupo.

Comparado con ACP, aquí las direcciones no se eligen por varianza total, sino
por capacidad de separación entre clases; esa diferencia de objetivo explica
por qué ambos métodos producen ejes distintos en general.

\begin{observacion}
Las variables canónicas quedan ordenadas por potencia discriminante
(autovalores decrecientes) y son incorreladas respecto a la métrica
intra-grupos.
\end{observacion}

% ══════════════════════════════════════════════════════════════════════════
\section{Evaluación del modelo: matriz de confusión}
% ══════════════════════════════════════════════════════════════════════════

\begin{definicion}[Matriz de confusión]
En clasificación binaria, una matriz de confusión recoge:
\begin{itemize}
  \item VP (verdaderos positivos), VN (verdaderos negativos),
        FP (falsos positivos), FN (falsos negativos).
  \item Según la convención habitual en aprendizaje automático:
        filas = clase real, columnas = clase predicha
        (algunos programas usan la convención traspuesta).
\end{itemize}
\end{definicion}

Métricas más usadas:
\begin{mdframed}[style=resultado]
\[
  \text{Exactitud}=\frac{VP+VN}{VP+VN+FP+FN},
\]
\[
  \text{Precisión (PPV)}=\frac{VP}{VP+FP},
  \qquad
  \text{Sensibilidad (Recall)}=\frac{VP}{VP+FN},
\]
\[
  \text{Especificidad}=\frac{VN}{VN+FP},
  \qquad
  \text{VPN}=\frac{VN}{VN+FN},
\]
\[
  F_1=\frac{2\,\text{Precisión}\cdot\text{Sensibilidad}}
            {\text{Precisión}+\text{Sensibilidad}}
      =\frac{2VP}{2VP+FP+FN}.
\]
\end{mdframed}

\begin{itemize}
  \item Tasa de falsos positivos: $\dfrac{FP}{FP+VN}=1-\text{Especificidad}$.
  \item Tasa de falsos negativos: $\dfrac{FN}{FN+VP}=1-\text{Sensibilidad}$.
\end{itemize}

\begin{nota}
Cuando las clases están desbalanceadas, conviene no evaluar solo con
exactitud; deben reportarse también sensibilidad, especificidad, precisión
y $F_1$, idealmente junto con validación cruzada estratificada.
\end{nota}

Este bloque cierra la metodología de clasificación: una buena regla teórica
debe validarse con métricas acordes al objetivo aplicado y al balance de clases.

% ══════════════════════════════════════════════════════════════════════════
\section{Estudio de las dimensiones necesarias}
% ══════════════════════════════════════════════════════════════════════════

No siempre se necesitan las $r=\min(p,G-1)$ dimensiones canónicas.
Para decidir cuántas conservar:

\begin{enumerate}
  \item Ordenar autovalores
        $\lambda_1\ge\lambda_2\ge\cdots\ge\lambda_r$ de $\mat{W}^{-1}\mat{B}$.
  \item Evaluar proporción acumulada de discriminación:
        \[
          \eta_q=\frac{\sum_{j=1}^{q}\lambda_j}{\sum_{j=1}^{r}\lambda_j}.
        \]
  \item Contrastar significación global o secuencial (p.\,ej. mediante
        estadísticos tipo Lambda de Wilks).
\end{enumerate}

Si normalizamos los autovectores para que
$\mathbf{v}_j'\mat{W}\mathbf{v}_j=1$, entonces:
\[
  \mathbf{v}_j'\mat{B}\mathbf{v}_j=\lambda_j,
\]
de forma que cada $\lambda_j$ cuantifica la separación adicional aportada
por la dimensión canónica $j$.

En consecuencia, la selección de dimensiones canónicas debe equilibrar potencia
discriminante, estabilidad muestral e interpretabilidad sustantiva.

% ══════════════════════════════════════════════════════════════════════════
\section{Algunas consideraciones prácticas}
% ══════════════════════════════════════════════════════════════════════════

\subsection{Hipótesis de partida}

Supuestos clásicos del ADL:
\begin{itemize}
  \item Normalidad multivariante dentro de cada grupo.
  \item Matriz de covarianzas común entre grupos
        ($\boldsymbol{\Sigma}_1=\cdots=\boldsymbol{\Sigma}_G$).
  \item Independencia muestral y ausencia de colinealidad extrema.
\end{itemize}

Estas hipótesis no deben leerse como formalidad: delimitan el régimen donde la
regla lineal es teóricamente óptima. Su comprobación condiciona la validez del
modelo y la interpretación de resultados.

Si la homogeneidad de covarianzas falla claramente, suele considerarse QDA.
Si la normalidad es débil, pueden funcionar bien transformaciones y/o métodos
no paramétricos.

\begin{itemize}
  \item Para evaluar homogeneidad de covarianzas se utiliza con frecuencia el
        test de \textbf{Box's M} (con cautela en muestras pequeñas).
  \item Si $p$ es grande respecto a $n$, conviene regularizar la matriz de
        covarianzas (\emph{shrinkage LDA}) para mejorar estabilidad y
        generalización.
\end{itemize}

\subsection{Tamaño muestral}

\begin{itemize}
  \item Regla práctica habitual: al menos 5 observaciones por predictor.
  \item Es recomendable que el grupo más pequeño tenga tamaño superior al
        número de predictores.
  \item Como guía operativa, se buscan al menos 20 observaciones por grupo
        para mayor estabilidad.
  \item Desbalances fuertes entre tamaños de grupo pueden sesgar la regla de
        clasificación hacia los grupos más frecuentes si no se corrigen las
        probabilidades a priori o los pesos de coste.
\end{itemize}

\begin{observacion}
Estas recomendaciones no son reglas rígidas; su función es controlar varianza
de estimación y evitar sobreajuste en escenarios de alta dimensión.
\end{observacion}

\begin{mdframed}[style=resultado]
Resumen operativo del Tema 4:
\begin{enumerate}
  \item Modelar $P(g\mid\mathbf{x})$ mediante Bayes y densidades de grupo.
  \item En LDA, usar funciones discriminantes lineales (o equivalentes por
        Mahalanobis).
  \item Para varios grupos, construir variables canónicas con
        $\mat{W}^{-1}\mat{B}$.
  \item Evaluar siempre la clasificación con matriz de confusión y métricas
        más allá de la exactitud.
\end{enumerate}
\end{mdframed}

\newpage
% ══════════════════════════════════════════════════════════════════════════
\section*{Formulario Final}
% ══════════════════════════════════════════════════════════════════════════
\addcontentsline{toc}{section}{Formulario Final}

\subsection*{Tema 1: Introducción y Álgebra Lineal}
\begin{mdframed}[style=resultado]
\textbf{Valor teórico (combinación lineal):}
\[
  z = w_1X_1 + w_2X_2 + \cdots + w_pX_p.
\]
\textbf{Matriz de datos:}
\[
  \mat{X}=[x_{ij}]_{n\times p}.
\]
\textbf{Producto escalar:}
\[
  \mathbf{x}'\mathbf{y}=\sum_i x_i y_i.
\]
\textbf{Norma euclídea:}
\[
  \|\mathbf{x}\|=\sqrt{\mathbf{x}'\mathbf{x}}.
\]
\textbf{Relación producto escalar--ángulo:}
\[
  \mathbf{x}'\mathbf{y}=\|\mathbf{x}\|\,\|\mathbf{y}\|\cos\theta.
\]
\textbf{Ecuación característica (autovalores):}
\[
  |\mat{A}-\lambda\mat{I}|=0.
\]
\textbf{Forma cuadrática:}
\[
  Q(\mathbf{x})=\mathbf{x}'\mat{A}\mathbf{x}
  =\sum_i a_{ii}x_i^2+\sum_{i<j}(a_{ij}+a_{ji})x_ix_j.
\]
\textbf{Inversa de matriz regular:}
\[
  \mat{A}^{-1}=\frac{1}{|\mat{A}|}\,\mat{A}^{\Lambda}.
\]
\end{mdframed}

\subsection*{Tema 2: Distribuciones Multivariantes}
\begin{mdframed}[style=resultado]
\textbf{Distribución conjunta (función de distribución):}
\[
  F(\mathbf{x}^0)=P(\mathbf{x}\le \mathbf{x}^0).
\]
\textbf{Paso de densidad a distribución:}
\[
  F(\mathbf{x}^0)=\int_{-\infty}^{\mathbf{x}^0} f(\mathbf{x})\,d\mathbf{x}.
\]
\textbf{Distribución condicionada:}
\[
  f(\mathbf{x}_1\mid \mathbf{x}_2^0)=
  \frac{f(\mathbf{x}_1,\mathbf{x}_2^0)}{f(\mathbf{x}_2^0)}.
\]
\textbf{Independencia:}
\[
  f(\mathbf{x}_1,\dots,\mathbf{x}_p)=\prod_{j=1}^p f(\mathbf{x}_j).
\]
\textbf{Vector de medias:}
\[
  \boldsymbol{\mu}=E(\mathbf{X}),\qquad
  \bar{\mathbf{x}}=\frac{1}{n}\sum_{i=1}^n\mathbf{x}_i.
\]
\textbf{Matriz de varianzas--covarianzas:}
\[
  \mat{S}= \frac{1}{n}\sum_{i=1}^n (\mathbf{x}_i-\bar{\mathbf{x}})(\mathbf{x}_i-\bar{\mathbf{x}})'.
\]
\textbf{Matriz de centrado y proyección:}
\[
  \tilde{\mathbf{X}}=\mat{P}\mathbf{X},\qquad
  \mat{P}=\mat{I}-\frac{1}{n}\mathbf{1}\mathbf{1}'.
\]
\textbf{Correlación lineal simple:}
\[
  r_{jk}=\frac{s_{jk}}{s_js_k}.
\]
\textbf{Correlación parcial:}
\[
  r_{jk.\,\text{resto}}=-\frac{s^{jk}}{\sqrt{s^{jj}s^{kk}}}.
\]
\textbf{Cambio de variable (jacobiano):}
\[
  f_y(\mathbf{y})=f_x(\mathbf{x})\left|\frac{d\mathbf{x}}{d\mathbf{y}}\right|.
\]
\textbf{Transformación lineal:}
\[
  \boldsymbol{\mu}_y=\mat{A}\boldsymbol{\mu}_x,\qquad
  \mat{S}_y=\mat{A}\mat{S}_x\mat{A}'.
\]
\textbf{Multinomial:}
\[
  P(\mathbf{n})=\frac{n!}{n_1!\cdots n_G!}\,p_1^{n_1}\cdots p_G^{n_G}.
\]
\textbf{Dirichlet:}
\[
  f(x_1,\dots,x_{G-1})=
  \frac{\Gamma(\alpha_0)}{\prod_{i=1}^G\Gamma(\alpha_i)}
  \prod_{i=1}^G x_i^{\alpha_i-1}.
\]
\textbf{Normal multivariante:}
\[
  f(\mathbf{x})=\frac{1}{(2\pi)^{p/2}|\boldsymbol{\Sigma}|^{1/2}}
  \exp\!\left\{-\frac12(\mathbf{x}-\boldsymbol{\mu})'\boldsymbol{\Sigma}^{-1}(\mathbf{x}-\boldsymbol{\mu})\right\}.
\]
\textbf{Normal condicionada (bloques):}
\[
  \mathbf{X}_1\mid \mathbf{X}_2=\mathbf{x}_2 \sim
  \mathcal{N}\!\left(
  \boldsymbol{\mu}_1+\boldsymbol{\Sigma}_{12}\boldsymbol{\Sigma}_{22}^{-1}(\mathbf{x}_2-\boldsymbol{\mu}_2),\;
  \boldsymbol{\Sigma}_{11}-\boldsymbol{\Sigma}_{12}\boldsymbol{\Sigma}_{22}^{-1}\boldsymbol{\Sigma}_{21}
  \right).
\]
\textbf{Distancia de Mahalanobis y ley $\chi^2$:}
\[
  D^2=(\mathbf{X}-\boldsymbol{\mu})'\boldsymbol{\Sigma}^{-1}(\mathbf{X}-\boldsymbol{\mu})\sim\chi^2(p).
\]
\end{mdframed}

\subsection*{Tema 3: Componentes Principales (ACP)}
\begin{mdframed}[style=resultado]
\textbf{Criterio geométrico básico:}
\[
  \min\sum_{i=1}^n r_i^2 \Longleftrightarrow \max\sum_{i=1}^n z_i^2.
\]
\textbf{Definición de componentes:}
\[
  \mat{Z}=\mathbf{X}\mat{A},\qquad \mat{A}'\mat{A}=\mat{I}.
\]
\textbf{Varianza del componente $j$:}
\[
  \operatorname{var}(\mathbf{z}_j)=\mathbf{a}_j'\mat{S}\mathbf{a}_j.
\]
\textbf{Problema de autovalores ACP:}
\[
  \mat{S}\mathbf{a}_j=\lambda_j\mathbf{a}_j,\qquad
  |\mat{S}-\lambda\mat{I}|=0.
\]
\textbf{Incorrelación entre componentes:}
\[
  \operatorname{cov}(\mathbf{z}_i,\mathbf{z}_j)=0\quad (i\neq j).
\]
\textbf{Conservación de varianza total:}
\[
  \tr(\mat{S})=\sum_{j=1}^p\lambda_j=\sum_{j=1}^p\operatorname{var}(z_j).
\]
\textbf{Proporción explicada por componente:}
\[
  c_h=\frac{\lambda_h}{\sum_{j=1}^p\lambda_j}.
\]
\textbf{Correlación componente--variable:}
\[
  r(z_i,x_j)=a_{ij}\frac{\sqrt{\lambda_i}}{s_j}.
\]
\textbf{Matriz de saturaciones (cargas factoriales):}
\[
  \mat{F}=\mat{A}\boldsymbol{\Lambda}^{1/2}.
\]
\textbf{ACP sobre correlaciones (análisis normado):}
\[
  \sum_{j=1}^p \lambda_j^{R}=\tr(\mat{R})=p,\qquad
  c_j=\frac{\lambda_j^{R}}{p}.
\]
\end{mdframed}

\subsection*{Tema 4: Análisis Discriminante (LDA/Fisher)}
\begin{mdframed}[style=resultado]
\textbf{Función discriminante lineal (dos clases):}
\[
  z_k=a+w_1x_{1k}+\cdots+w_px_{pk}.
\]
\textbf{Posterior bayesiana (dos clases):}
\[
  P(\mathcal{P}_i\mid\mathbf{x}_0)=
  \frac{\pi_i f_i(\mathbf{x}_0)}
       {\pi_1 f_1(\mathbf{x}_0)+\pi_2 f_2(\mathbf{x}_0)}.
\]
\textbf{Regla con costes de clasificación:}
\[
  c(1|2)\,\pi_2 f_2(\mathbf{x}_0) > c(2|1)\,\pi_1 f_1(\mathbf{x}_0)
  \Longrightarrow \text{clasificar en }\mathcal{P}_2.
\]
\textbf{Discriminantes lineales gaussianos (covarianza común):}
\[
  \delta_i(\mathbf{x})=
  \mathbf{x}'\boldsymbol{\Sigma}^{-1}\boldsymbol{\mu}_i
  -\frac12\boldsymbol{\mu}_i'\boldsymbol{\Sigma}^{-1}\boldsymbol{\mu}_i
  +\log\pi_i.
\]
\textbf{Distancia de Mahalanobis por grupo:}
\[
  D_i^2(\mathbf{x})=(\mathbf{x}-\boldsymbol{\mu}_i)'\boldsymbol{\Sigma}^{-1}(\mathbf{x}-\boldsymbol{\mu}_i).
\]
\textbf{Posterior en función de distancias (dos grupos):}
\[
  P(\mathcal{P}_1\mid\mathbf{x})=
  \frac{1}{1+\dfrac{\pi_2}{\pi_1}
  \exp\!\left[-\frac12\left(D_2^2(\mathbf{x})-D_1^2(\mathbf{x})\right)\right]}.
\]
\textbf{Estimadores muestrales en LDA:}
\[
  \hat{\pi}_g=\frac{n_g}{n},\qquad
  \hat{\boldsymbol{\mu}}_g=\frac{1}{n_g}\sum_{i:y_i=g}\mathbf{x}_i,\qquad
  \hat{\boldsymbol{\Sigma}}=\frac{1}{n-G}\sum_{g=1}^{G}\sum_{i:y_i=g}
  (\mathbf{x}_i-\hat{\boldsymbol{\mu}}_g)(\mathbf{x}_i-\hat{\boldsymbol{\mu}}_g)'.
\]
\textbf{Error teórico simétrico (dos gaussianas):}
\[
  \Delta^2=(\boldsymbol{\mu}_1-\boldsymbol{\mu}_2)'\boldsymbol{\Sigma}^{-1}(\boldsymbol{\mu}_1-\boldsymbol{\mu}_2),
  \qquad
  P(2|1)=P(1|2)=1-\Phi\!\left(\frac{\Delta}{2}\right).
\]
\textbf{Región de decisión multiclase:}
\[
  \mathcal{A}_g=\{\mathbf{x}\in\R^p:\pi_g f_g(\mathbf{x})>\pi_i f_i(\mathbf{x}),\ \forall i\neq g\}.
\]
\textbf{Frontera lineal entre clases $i,j$:}
\[
  (\boldsymbol{\mu}_i-\boldsymbol{\mu}_j)'\boldsymbol{\Sigma}^{-1}\mathbf{x}
  -\frac12\!\left(\boldsymbol{\mu}_i'\boldsymbol{\Sigma}^{-1}\boldsymbol{\mu}_i-\boldsymbol{\mu}_j'\boldsymbol{\Sigma}^{-1}\boldsymbol{\mu}_j\right)
  +\log\frac{\pi_i}{\pi_j}=0.
\]
\textbf{Criterio de Fisher (dos grupos):}
\[
  J(\boldsymbol{\alpha})=
  \frac{\left[\boldsymbol{\alpha}'(\bar{\mathbf{x}}_2-\bar{\mathbf{x}}_1)\right]^2}
       {\boldsymbol{\alpha}'\mat{S}_w\boldsymbol{\alpha}},
  \qquad
  \boldsymbol{\alpha}\propto\mat{S}_w^{-1}(\bar{\mathbf{x}}_2-\bar{\mathbf{x}}_1).
\]
\textbf{Dispersión intra/intergrupo y problema generalizado:}
\[
  \mat{W}=\sum_{g=1}^{G}\sum_{i=1}^{n_g}(\mathbf{x}_{ig}-\bar{\mathbf{x}}_g)(\mathbf{x}_{ig}-\bar{\mathbf{x}}_g)',\quad
  \mat{B}=\sum_{g=1}^{G} n_g(\bar{\mathbf{x}}_g-\bar{\mathbf{x}})(\bar{\mathbf{x}}_g-\bar{\mathbf{x}})',
\]
\[
  \mat{B}\mathbf{u}=\lambda\mat{W}\mathbf{u}
  \Longleftrightarrow
  \mat{W}^{-1}\mat{B}\mathbf{u}=\lambda\mathbf{u}.
\]
\textbf{Variables canónicas y número máximo:}
\[
  \mathbf{z}=\mat{U}_r'\mathbf{x},\qquad r=\min(p,G-1).
\]
\textbf{Proporción acumulada de discriminación:}
\[
  \eta_q=\frac{\sum_{j=1}^{q}\lambda_j}{\sum_{j=1}^{r}\lambda_j}.
\]
\textbf{Métricas de clasificación (matriz de confusión):}
\begin{align}
  \text{Exactitud} &= \frac{VP+VN}{VP+VN+FP+FN}, \\
  \text{Precisión} &= \frac{VP}{VP+FP}, \\
  \text{Sensibilidad} &= \frac{VP}{VP+FN}, \\
  \text{Especificidad} &= \frac{VN}{VN+FP}, \\
  F_1 &= \frac{2VP}{2VP+FP+FN}.
\end{align}
\end{mdframed}

\end{document}
